%% 子ども用『物理を理解するための数学　ガイドブック　はじめに読む本』親ファイル（第16回編集委員会、組 G-A）
%% vol=0（巻の色 bkVol0＝紺灰色）。ルビつき（第2・3巻と同じ）。表紙に「子ども用」の札（丸い形）。
%% 表1・表2・とびら（p.1）・もくじ（p.2）・はじめに（p.3）・本文 guide-p1（p.4〜、節は偶数ページから）・奥付・表4。
%% もくじは \pageref で各節のラベル kga-sec1〜kga-sec9 を拾う。共通の部品は guide-kga.tex。
%% 組み方：cd tex && lualatex guide-kids.tex（2回。もくじのページ番号は2回目で入る）
\documentclass[12pt]{ltjsarticle}
\def\bkfontpath{fonts/}
\usepackage[vol=0]{bk-b5}
\usepackage{bk-matter}
\usepackage{bk-honbun}
\usepackage{bk-volume}
\usepackage{bk-pict1}
%% bk-motifs.tex  表紙の図柄（各巻の数学の考え方を絵にしたもの）
%% 描く範囲：x 0〜150, y 0〜106（mm）。色 mA＝主（白か濃い色）, mB＝従（淡い色か白）, mC＝影
\makeatletter
\tikzset{mot/.style={line cap=round, line join=round}}

%% 第1巻：10のまとまり・ものさし・時計
\newcommand\bk@apple[2]{\bkapple{(#1,#2)}{4.3}{bkAppleRed}}% リンゴ（へた・葉つき）
\newcommand\bkMotifi{%
  % 10のわく：白いマスにリンゴを1こずつ
  \foreach \i in {0,...,4} \foreach \j in {0,1} {
    \fill[white, rounded corners=1.6mm] (2.6+\i*14.4,48.6+\j*14.4) rectangle ++(13.2,13.2);
    \bk@apple{9.2+\i*14.4}{54.6+\j*14.4}}
  % のこりの3こは、白いおさらの上に
  \bkplate{(101,56)}{38}
  \foreach \i in {0,1,2} \bk@apple{89+\i*12}{61.5}
  \node[mA, font=\sffamily\bfseries\fontsize{22}{22}\selectfont] at (38,39) {10};
  \node[mA, font=\sffamily\bfseries\fontsize{22}{22}\selectfont] at (101,42) {3};
  % 時計（3時）
  \fill[mB] (130,88) circle (15.5);
  \draw[mA, line width=1.6pt] (130,88) circle (15.5);
  \foreach \a in {0,30,...,330} \draw[mC, line width=1pt] (130,88) ++(\a:12.6) -- ++(\a:2.2);
  \draw[mC, line width=2pt, line cap=round] (130,88) -- ++(0:8);
  \draw[mC, line width=1.4pt, line cap=round] (130,88) -- ++(90:11.5);
  \fill[mC] (130,88) circle (1.3);
  % ものさし（0〜10）
  \fill[mB, rounded corners=1.6mm] (2,6) rectangle (122,24);
  \foreach \k in {0,...,20} {\pgfmathsetmacro\bkx{7+\k*5.5}
    \ifodd\k \draw[mC, line width=0.7pt] (\bkx,24) -- (\bkx,20.5);
    \else \draw[mC, line width=0.9pt] (\bkx,24) -- (\bkx,18.5);\fi}
  \foreach \k in {0,...,10} \node[mC, font=\sffamily\bfseries\fontsize{7.5}{8}\selectfont] at (7+\k*11,13.5) {\k};
}

%% 第2巻：落ちるボール（同じ時間ごと）・倍々・単位・数直線
\newcommand\bkMotifii{%
  \draw[mA, line width=1.4pt] (5,2) -- (5,104);
  \foreach \y in {2,12.2,...,104} \draw[mA, line width=0.9pt] (5,\y) -- (8.5,\y);
  \foreach \t/\y in {0/100, 1/93.6, 2/74.5, 3/42.7, 4/-1.9} {
    \fill[mA] (16,{\y+3}) circle (3.0);
    \node[anchor=west, mA, font=\sffamily\bfseries\fontsize{8.5}{10}\selectfont] at (20.5,{\y+3}) {\t\R{秒}{びょう}};}
  % 単位のふだ
  \foreach \x/\t in {92/m, 112/kg, 132/s} {
    \fill[mB, rounded corners=2mm] (\x,88) rectangle ++(17,13);
    \node[mC, font=\rmfamily\bfseries\fontsize{17}{17}\selectfont] at (\x+8.5,94.5) {\t};}
  % 倍々のブロック（1, 2, 4, 8）
  \begin{scope}[shift={(41,44)}]
    \foreach \bx/\bw/\bh/\bn in {0/1/1/1, 19.5/2/1/2, 47.5/2/2/4, 75.5/4/2/8} {
      \fill[mA] (\bx,0) rectangle ++(\bw*8.5,\bh*8.5);
      \draw[mC, line width=0.7pt, shift={(\bx,0)}] (0,0) grid[step=8.5] ({\bw*8.5},{\bh*8.5});
      \node[mA, font=\sffamily\bfseries\fontsize{16}{16}\selectfont] at (\bx+\bw*4.25,\bh*8.5+6.5) {\bn};}
    \foreach \xm in {14, 42, 70} \node[mB, font=\sffamily\bfseries\fontsize{11}{11}\selectfont] at (\xm,4.25) {×2};
  \end{scope}
  % 数直線（−5〜5）
  \draw[mA, line width=1.6pt, -{Stealth[length=8pt]}] (40,14) -- (150,14);
  \foreach \k in {-5,...,5} {\pgfmathsetmacro\bkx{93+\k*10}
    \draw[mA, line width=1.1pt] (\bkx,11.5) -- (\bkx,16.5);
    \node[mA, font=\sffamily\bfseries\fontsize{9.5}{9.5}\selectfont] at (\bkx,6.5) {\ifnum\k<0 −\number-\k\else\k\fi};}% 第7回：数字をゴシックに（マイナスは記号の字）
  \draw[mB, line width=3.4pt, -{Stealth[length=9pt]}] (93,24) -- (123,24);
  \draw[mB, line width=3.4pt, -{Stealth[length=9pt]}] (93,24) -- (73,24);
  \fill[mA] (93,24) circle (2.4);
}

%% 第3巻：データに直線・投げたボールの道・ヒストグラム
\newcommand\bkMotifiii{%
  \draw[mA, line width=1.3pt, -{Stealth[length=7pt]}] (4,6) -- (98,6);
  \draw[mA, line width=1.3pt, -{Stealth[length=7pt]}] (6,4) -- (6,104);
  \draw[mB, line width=4pt] (10,12) -- (94,86);
  \foreach \x/\e in {15/3, 24/-4, 33/2.5, 42/-2, 51/4, 60/-3, 69/1.5, 78/-1.5, 87/2.5}
    \fill[mA] (\x,{12+(\x-10)*74/84+\e}) circle (2.3);
  \draw[mA, line width=1.3pt, dash pattern=on 4pt off 3pt] plot[domain=104:148, samples=40] (\x,{8+74*(1-((\x-126)/22)^2)});
  \foreach \x in {104,111.33,118.67,126,133.33,140.67,148} \fill[mA] (\x,{8+74*(1-((\x-126)/22)^2)}) circle (2.4);
  \foreach \h [count=\i] in {2,5,10,15,18,15,10,5,2}
    \fill[mB] (99+\i*5.4,88) rectangle ++(4.8,\h);
  \draw[mA, line width=1pt] (103,88) -- (150,88);
}

%% 第4巻：円運動の影＝波・ベクトルの合成
\newcommand\bkMotifiv{%
  \def\bkcy{64}
  \draw[mB, line width=0.9pt] (2,\bkcy) -- (66,\bkcy) (34,\bkcy-32) -- (34,\bkcy+32);
  \draw[mA, line width=1.6pt] (34,\bkcy) circle (28);
  \fill[mB] (34,\bkcy) -- (42,\bkcy) arc (0:50:8) -- cycle;
  \draw[mA, line width=2pt] (34,\bkcy) -- ++(50:28);
  \draw[mC, line width=1.4pt, dash pattern=on 3pt off 2pt] (52,{\bkcy+28*sin(50)}) -- (52,\bkcy);
  \draw[mC, line width=3pt] (34,\bkcy) -- (52,\bkcy);
  \fill[mA] (52,{\bkcy+28*sin(50)}) circle (2.4);
  \draw[mA, line width=1pt, dash pattern=on 3pt off 2pt] (54,{\bkcy+28*sin(50)}) -- (72,{\bkcy+28*sin(50)});
  \draw[mB, line width=0.9pt] (72,\bkcy) -- (150,\bkcy);
  \draw[mA, line width=2.2pt] plot[domain=72:150, samples=120, smooth] (\x,{\bkcy+28*sin(50+(\x-72)*5.142857)});
  % ベクトル
  \begin{scope}[shift={(92,4)}]
    \draw[mB, line width=1pt, dash pattern=on 3pt off 2pt] (32,7) -- (44,22) -- (12,15);
    \draw[mA, line width=2pt, -{Stealth[length=8pt]}] (0,0) -- (32,7);
    \draw[mA, line width=2pt, -{Stealth[length=8pt]}] (0,0) -- (12,15);
    \draw[mC, line width=2.6pt, -{Stealth[length=9pt]}] (0,0) -- (44,22);
  \end{scope}
}

%% 第5巻：接線（微分）と細い長方形（積分）
\newcommand\bkMotifv{%
  \draw[mA, line width=1.3pt, -{Stealth[length=7pt]}] (4,6) -- (150,6);
  \draw[mA, line width=1.3pt, -{Stealth[length=7pt]}] (6,4) -- (6,104);
  \foreach \k in {0,...,7} {\pgfmathsetmacro\bkx{16+\k*10}
    \pgfmathsetmacro\bkh{16+0.35*((\bkx-6)/10)^2+9*sin(3*(\bkx-6))}
    \fill[mB] (\bkx,6) rectangle (\bkx+10,\bkh);
    \draw[mC, line width=0.6pt] (\bkx,6) rectangle (\bkx+10,\bkh);}
  \draw[mA, line width=2.4pt] plot[domain=8:146, samples=120, smooth] (\x,{16+0.35*((\x-6)/10)^2+9*sin(3*(\x-6))});
  \pgfmathsetmacro\bky{16+0.35*11.4^2+9*sin(3*114)}
  \pgfmathsetmacro\bks{0.007*114+9*3*pi/180*cos(3*114)}
  \draw[mA, line width=1.2pt, dash pattern=on 3pt off 2pt] (120,\bky) -- (134,\bky) -- (134,\bky+14*\bks);% 第6回：緑地で見えるように白に
  \draw[mA, line width=1.8pt] (100,{\bky-20*\bks}) -- (146,{\bky+26*\bks});
  \fill[mA] (120,\bky) circle (2.6);
  \node[mA, anchor=north, font=\sffamily\bfseries\fontsize{10}{10}\selectfont] at (127,\bky-1.8) {$\Delta x$};% 第6回：横の辺の中央の下
  \node[mA, anchor=west, font=\sffamily\bfseries\fontsize{10}{10}\selectfont] at (135,\bky+8) {$\Delta y$};
}

%% 第6巻：ゆれながら止まる（減衰振動）・回る矢印（複素数）
\newcommand\bkMotifvi{%
  \draw[mB, line width=0.9pt] (6,56) -- (150,56);
  \draw[mA, line width=1.3pt, -{Stealth[length=7pt]}] (6,6) -- (6,106);
  \draw[mB, line width=1.2pt, dash pattern=on 4pt off 3pt] plot[domain=8:148, samples=80] (\x,{56+40*exp(-0.022*(\x-8))});
  \draw[mB, line width=1.2pt, dash pattern=on 4pt off 3pt] plot[domain=8:148, samples=80] (\x,{56-40*exp(-0.022*(\x-8))});
  \draw[mA, line width=2.4pt] plot[domain=8:148, samples=200, smooth] (\x,{56+40*exp(-0.022*(\x-8))*cos(9.17*(\x-8))});
  \begin{scope}[shift={(132,90)}]
    \draw[mB, line width=0.8pt] (-15,0) -- (15,0) (0,-14) -- (0,15);
    \draw[mA, line width=1.2pt] (0,0) circle (11);
    \draw[mA, line width=2pt, -{Stealth[length=7pt]}] (0,0) -- (40:11);
    \draw[mB, line width=1.4pt, -{Stealth[length=6pt]}] (55:13.5) arc (55:120:13.5);
  \end{scope}
}

%% 第7巻：行列で変わる格子
\newcommand\bkMotifvii{%
  \begin{scope}
    \clip (0,0) rectangle (150,106);
    \begin{scope}[shift={(72,50)}]
      \foreach \i in {-7,...,7} {
        \draw[mB, line width=0.5pt] (\i*12,-60) -- (\i*12,60);
        \draw[mB, line width=0.5pt] (-90,\i*12) -- (90,\i*12);}
      \begin{scope}[cm={1,0.28,0.48,1,(0,0)}]
        \foreach \i in {-8,...,8} {
          \draw[mA, line width=0.9pt] (\i*12,-80) -- (\i*12,80);
          \draw[mA, line width=0.9pt] (-110,\i*12) -- (110,\i*12);}
        \fill[mC] (0,0) rectangle (12,12);
        \draw[mA, line width=2.4pt, -{Stealth[length=9pt]}] (0,0) -- (24,0);
        \draw[mA, line width=2.4pt, -{Stealth[length=9pt]}] (0,0) -- (0,24);
      \end{scope}
      \fill[mA] (0,0) circle (2);
    \end{scope}
  \end{scope}
}

%% 第8巻：わき出し（div）と うず（rot）
\newcommand\bkMotifviii{%
  % うすい等高線と、いちばん急な向き（grad）
  \foreach \r in {6,12,...,30} \draw[mB, line width=0.5pt] (75,94) ellipse ({\r*1.5} and {\r*0.3});
  % わき出し
  \begin{scope}[shift={(36,50)}]
    \fill[mB] (0,0) circle (5);\fill[mA] (0,0) circle (2.4);
    \foreach \a in {0,30,...,330} \draw[mA, line width=1.6pt, -{Stealth[length=5pt]}] (\a:9) -- (\a:16);
    \foreach \a in {15,45,...,345} \draw[mA, line width=1.3pt, -{Stealth[length=5pt]}] (\a:21) -- (\a:28);
    \node[mB, font=\sffamily\bfseries\fontsize{14}{14}\selectfont] at (0,-38) {div};
  \end{scope}
  % うず
  \begin{scope}[shift={(112,50)}]
    \fill[mB] (0,0) circle (2.4);
    \foreach \a in {0,45,...,315} \draw[mA, line width=1.6pt, -{Stealth[length=5pt]}] (\a:11) ++(\a-90:2.6) -- ++(\a+90:5.2);
    \foreach \a in {15,45,...,345} \draw[mA, line width=1.3pt, -{Stealth[length=5pt]}] (\a:24.5) ++(\a-90:3.6) -- ++(\a+90:7.2);
    \node[mB, font=\sffamily\bfseries\fontsize{14}{14}\selectfont] at (0,-38) {rot};
  \end{scope}
}

%% 第9巻：sinをたして四角い波・スペクトル
\newcommand\bkMotifix{%
  \def\bkmid{72}\def\bkamp{15}
  \draw[mB, line width=1.2pt, dash pattern=on 4pt off 3pt] (8,\bkmid-\bkamp) -- (8,\bkmid+\bkamp) -- (43,\bkmid+\bkamp) -- (43,\bkmid-\bkamp) -- (78,\bkmid-\bkamp) -- (78,\bkmid+\bkamp) -- (113,\bkmid+\bkamp) -- (113,\bkmid-\bkamp) -- (148,\bkmid-\bkamp) -- (148,\bkmid);
  \draw[mB, line width=1.1pt] plot[domain=8:148, samples=120, smooth] (\x,{\bkmid+\bkamp*4/pi*sin((\x-8)*5.142857)});
  \draw[mA, line width=2.2pt] plot[domain=8:148, samples=320, smooth] (\x,{\bkmid+\bkamp*4/pi*(sin((\x-8)*5.142857)+sin(3*(\x-8)*5.142857)/3+sin(5*(\x-8)*5.142857)/5+sin(7*(\x-8)*5.142857)/7+sin(9*(\x-8)*5.142857)/9)});
  \draw[mA, line width=1.3pt] (8,8) -- (148,8);
  \foreach \k [count=\i] in {1,3,5,7,9,11} {
    \fill[mA] ({4+\i*18},8) rectangle ++(9,{36/\k});
    \node[mB, font=\sffamily\bfseries\fontsize{8}{8}\selectfont] at ({8.5+\i*18},4) {\k};}
}
\makeatother

%% 自動生成（gen.py）。直接編集しないこと
\makeatletter
\bkdefvol{1}{かぞえる・はかる}{}{かずで くらべる、かずで つたえる}{小学1年〜3年}{1}{1}
\bkdefvol{2}{\R{数}{かず}と\R{量}{りょう}のことば}{}{「5」ではなく「5\,m」。\R{数}{かず}と\R{量}{りょう}のことばを\R{身}{み}につけよう}{小学4年〜中学1年}{1}{2}
\bkdefvol{3}{\R{関数}{かん|すう}とデータ}{きまりを\R{見}{み}つける}{\R{表}{ひょう}・グラフ・\R{式}{しき}で、\R{変化}{へん|か}のきまりをつかもう}{中学1年〜3年}{1}{3}
\bkdefvol{4}{\R{形}{かたち}と\R{向}{む}き}{\R{図形}{ず|けい}・\R{三角関数}{さん|かく|かん|すう}・ベクトル}{長さと角度で、形と向きを数にしよう}{中学2年〜高校2年}{1}{5}
\bkdefvol{5}{\R{変化}{へん|か}をつかむ}{\R{微分}{び|ぶん}と\R{積分}{せき|ぶん}}{変わる速さと、\mbox{たまった量}を計算しよう}{高校1年〜2年}{2}{4}
\bkdefvol{6}{\R{未来}{み|らい}を\R{計算}{けい|さん}する}{\R{微分方程式}{び|ぶん|ほう|てい|しき}と\R{複素数}{ふく|そ|すう}}{「次の瞬間のルール」から、未来を計算しよう}{高校2年〜3年}{2}{5}
\bkdefvol{7}{まとめて\R{計算}{けい|さん}する}{\R{行列}{ぎょう|れつ}}{たくさんの数を、表のまま計算しよう}{高校3年（発展）〜大学1年}{3}{5}
\bkdefvol{8}{\R{空間}{くう|かん}に\R{広}{ひろ}がる\R{量}{りょう}}{ベクトル\R{解析}{かい|せき}}{空間に広がる量を、矢印で読みとこう}{大学1年〜2年}{3}{5}
\bkdefvol{9}{\R{波}{なみ}と\R{変化}{へん|か}を\R{読}{よ}みとく}{フーリエ\R{解析}{かい|せき}とラプラス\R{変換}{へん|かん}}{波を成分に分け、微分方程式をかけ算で解こう}{大学2年}{3}{5}
\@namedef{bk@rname@1}{\R{変化}{へん|か}}\@namedef{bk@rsub@1}{ひき\R{算}{ざん}→\R{傾}{かたむ}き→\R{微分}{び|ぶん}}\@namedef{bk@rend@1}{\R{流体}{りゅう|たい}・\R{熱}{ねつ}}
\@namedef{bk@rname@2}{たし\R{合}{あ}わせ}\@namedef{bk@rsub@2}{かけ\R{算}{ざん}→\R{面積}{めん|せき}→\R{積分}{せき|ぶん}}\@namedef{bk@rend@2}{\R{電磁気}{でん|じ|き}}
\@namedef{bk@rname@3}{\R{倍々}{ばい|ばい}}\@namedef{bk@rsub@3}{\R{累乗}{るい|じょう}→\R{指数}{し|すう}→$e$}\@namedef{bk@rend@3}{\R{制御}{せい|ぎょ}}
\@namedef{bk@rname@4}{\R{回転}{かい|てん}と\R{波}{なみ}}\@namedef{bk@rsub@4}{\R{角}{かく}→sin→\R{複素数}{ふく|そ|すう}}\@namedef{bk@rend@4}{\R{音}{おと}・\R{交流}{こう|りゅう}}
\@namedef{bk@rname@5}{\R{表}{ひょう}}\@namedef{bk@rsub@5}{\R{九九}{く|く}の\R{表}{ひょう}→\R{行列}{ぎょう|れつ}}\@namedef{bk@rend@5}{\R{回路網}{かい|ろ|もう}・\R{制御}{せい|ぎょ}}
\@namedef{bk@rname@6}{\R{矢印}{や|じるし}}\@namedef{bk@rsub@6}{\R{向}{む}き→ベクトル→\R{場}{ば}}\@namedef{bk@rend@6}{ロボット・\\\R{電磁気}{でん|じ|き}}
\@namedef{bk@rname@7}{はかる}\@namedef{bk@rsub@7}{\R{目}{め}もり→\R{誤差}{ご|さ}→\R{統計}{とう|けい}}\@namedef{bk@rend@7}{\R{計測}{けい|そく}・\R{診断}{しん|だん}}
\makeatother

%% 自動生成（gen.py）
\makeatletter
\def\bk@npart{3}
\def\bk@routepage{4}
\bkdefpart{1}{第1部}{フーリエ解析}{（9-1〜9-4）}
\bkdefpart{2}{第2部}{ラプラス変換}{（9-5〜9-9）}
\bkdefpart{3}{第3部}{行列と微分方程式}{（9-10〜9-12）}
\@namedef{bk@rhere@1}{（この\R{巻}{かん}には\R{駅}{えき}がありません）}
\@namedef{bk@rhere@2}{（この\R{巻}{かん}には\R{駅}{えき}がありません）}
\@namedef{bk@rhere@3}{9-5、9-6、9-7、9-8、9-9}
\@namedef{bk@rhere@4}{9-1、9-2、9-3}
\@namedef{bk@rhere@5}{9-10、9-11}
\@namedef{bk@rhere@6}{（この\R{巻}{かん}には\R{駅}{えき}がありません）}
\@namedef{bk@rhere@7}{9-4、9-12}
\def\bk@stations{\bkstation{1}{1}{1-5}\bkstation{1}{2}{2-3}\bkstation{1}{3}{3-5}\bkstation{1}{5}{5-1}\bkstation{1}{6}{6-1}\bkstation{1}{8}{8-4}\bkstation{2}{1}{1-21}\bkstation{2}{2}{2-3}\bkstation{2}{4}{4-5}\bkstation{2}{5}{5-10}\bkstation{2}{6}{6-3}\bkstation{2}{8}{8-11}\bkstation{3}{2}{2-7}\bkstation{3}{3}{3-11}\bkstation{3}{5}{5-14}\bkstation{3}{6}{6-4}\bkstation{3}{9}{9-5}\bkstation{4}{1}{1-20}\bkstation{4}{4}{4-8}\bkstation{4}{5}{5-15}\bkstation{4}{6}{6-10}\bkstation{4}{9}{9-1}\bkstation{5}{1}{1-13}\bkstation{5}{2}{2-18}\bkstation{5}{3}{3-6}\bkstation{5}{7}{7-1}\bkstation{5}{9}{9-10}\bkstation{6}{2}{2-2}\bkstation{6}{3}{3-2}\bkstation{6}{4}{4-10}\bkstation{6}{7}{7-7}\bkstation{6}{8}{8-5}\bkstation{7}{1}{1-11}\bkstation{7}{2}{2-6}\bkstation{7}{3}{3-14}\bkstation{7}{9}{9-4}}
\@namedef{bk@rtop@1}{182}
\@namedef{bk@rtop@2}{182}
\@namedef{bk@rtop@3}{165}
\@namedef{bk@rtop@4}{182}
\@namedef{bk@rtop@5}{182}
\@namedef{bk@rtop@6}{165}
\@namedef{bk@rtop@7}{182}
\makeatother

\hypersetup{pdftitle={物理を理解するための数学　ガイドブック　はじめに読む本（子ども用）}, pdfauthor={「物理を理解するための数学」編集委員会}}
\newcommand\kgaBookTitle{ガイドブック}
\newcommand\kgaBookSub{はじめに\R{読}{よ}む\R{本}{ほん}}
\newcommand\kgaCatch{\R{全}{ぜん}9\R{巻}{かん}の\R{地図}{ち|ず}と、\R{歩}{ある}き\R{方}{かた}}
\newcommand\kgaBandLabel{ガイドブック\hspace{0.6\zw}はじめに\R{読}{よ}む\R{本}{ほん}}
%% guide-kga.tex　ガイドブック2冊（子ども用 guide-kids.tex・大人用 guide-adult.tex）で共通の部品（組 G-A、頭文字 kga）
%% 第16回編集委員会（10月30日 10:14 の依頼者の指示で2冊に分けた）。親ファイルの前付けと奥付・子ども用の本文で使う。
\makeatletter
%% 色：kgaPA＝子ども用の札の色（れんが色。白い地で約5.9:1）、kgaPB＝大人用の札の色（巻の色 bkVol0）
%% 札は色だけでなく、字（子ども用／大人用）と形（丸／四角）でも区別する
\definecolor{kgaPA}{cmyk}{0.05,0.65,0.95,0.28}
\colorlet{kgaPB}{bkVol0}
\newcommand\kgaNoRuby{\def\R##1##2{##1}}% ルビを外す
%% 本ごとに親ファイルで決めるもの（書名・副題・キャッチ・帯の小さな字・偶数ページの柱）
\providecommand\kgaBookTitle{ガイドブック}
\providecommand\kgaBookSub{}
\providecommand\kgaCatch{}
\providecommand\kgaBandLabel{}
%% ラベルのページ番号（まだ組んでいない所は「—」）  \kgaPage{ラベル}
\newcommand\kgaPage[1]{\@ifundefined{r@#1}{—}{\pageref{#1}}}

%% --------------------------------------------------------------------
%% 9冊の本のたなの絵  \kgaShelf{x}{y}{倍率}{棚板の色}：本の高さはページ数に比例（400〜640ページ）
%% --------------------------------------------------------------------
\newcommand\kgaShelf[4]{%
  \begin{scope}[shift={(#1,#2)}, scale=#3]
    \foreach \v/\p in {1/400,2/456,3/564,4/560,5/492,6/600,7/492,8/632,9/640} {
      \pgfmathsetmacro\kga@bx{(\v-1)*16.2}\pgfmathsetmacro\kga@bh{\p*0.13}
      \fill[bkVol\v, rounded corners=0.8mm] (\kga@bx,0) rectangle ++(14.4,\kga@bh);
      \draw[bkBandOn\v, line width=0.5pt, opacity=0.55] (\kga@bx+1.6,\kga@bh-6) -- ++(11.2,0) (\kga@bx+1.6,6) -- ++(11.2,0);
      \node[text=bkBandOn\v, font=\sffamily\bfseries\fontsize{15}{15}\selectfont] at (\kga@bx+7.2,\kga@bh-14) {\v};
    }
    \fill[#4, rounded corners=0.6mm] (-3,-3.2) rectangle (148.2,0);
  \end{scope}}

%% --------------------------------------------------------------------
%% 表紙の札  \kgaBadge{kids|adult}{中心x}{中心y}{大きさ（1＝表紙）}
%%   子ども用：白い丸に、れんが色の太い輪と「子ども用」。大人用：白い四角に、紺灰色の二重の枠と「大人用」
%% --------------------------------------------------------------------
\newcommand\kgaBadge[4]{%
  \begin{scope}[shift={(#2,#3)}, scale=#4, transform shape]
  \def\kga@b{#1}\def\kga@kids{kids}%
  \ifx\kga@b\kga@kids
    \fill[white] (0,0) circle (14);
    \draw[kgaPA, line width=2.4pt] (0,0) circle (12.6);
    \node[text=kgaPA, font=\sffamily\bfseries\fontsize{15}{17}\selectfont, inner sep=0pt] at (0,0) {\R{子}{こ}ども\R{用}{よう}};
  \else
    \fill[white] (-17,-11) rectangle (17,11);
    \draw[kgaPB, line width=1.6pt] (-15.4,-9.4) rectangle (15.4,9.4);
    \draw[kgaPB, line width=0.6pt] (-13.6,-7.6) rectangle (13.6,7.6);
    \node[text=kgaPB, font=\sffamily\bfseries\fontsize{14}{16}\selectfont, inner sep=0pt] at (0,0) {大人用};
  \fi
  \end{scope}}

%% --------------------------------------------------------------------
%% 表1（表紙）：\bkFrontArt と同じ配置。巻の番号の札のかわりに「子ども用」「大人用」の札
%%   \kgaFrontArt{kids|adult}
%% --------------------------------------------------------------------
\newcommand\kgaFrontArt[1]{%
  \fill[bkVol] (-3,64) rectangle (185,260);
  \fill[white] (-3,-3) rectangle (185,64);
  \node[anchor=base west, inner sep=0pt, text=bkOn, font=\sffamily\bfseries\fontsize{12.5}{15}\selectfont] at (16,242) {\bkSeries};
  \draw[bkOn, line width=0.6pt] (16,237.5) -- (112,237.5);
  \kgaBadge{#1}{147}{236}{1}
  \node[anchor=base west, inner sep=0pt, text=bkOn, font=\bk@titlefont] at (15,206) {\kgaBookTitle};
  \node[anchor=base west, inner sep=0pt, text=bkOn, font=\sffamily\bfseries\fontsize{17}{20}\selectfont] at (16,192) {\kgaBookSub};
  \node[anchor=base west, inner sep=0pt, text=bkOn, font=\sffamily\fontsize{11.5}{14}\selectfont] at (16,181.5) {\kgaCatch};
  \foreach \r in {1,...,7} {\node[inner sep=0pt] at ({20+(\r-1)*11.5},164) {\bkrouteicon[7.5mm]{\r}};}
  \kgaShelf{18}{74}{1}{white}
  \bk@lineup{16}{36}
  \node[anchor=base, text=bkGray, font=\sffamily\fontsize{8}{10}\selectfont] at (91,16) {\R{全}{ぜん}9\R{巻}{かん}　\bkSeries};
  \node[anchor=base, text=bkGray, font=\sffamily\fontsize{7.5}{9}\selectfont] at (91,8.5) {\bkEditor\hspace{1pt}\R{編}{へん}};
}

%% --------------------------------------------------------------------
%% 表4（裏表紙）  \kgaBackArt{見出し}{中身（TikZ。y=200 から下へ）}{ふきだしの文}
%% --------------------------------------------------------------------
\newcommand\kgaBackArt[3]{%
  \fill[bkVol!9] (-3,-3) rectangle (185,260);
  \fill[bkVol] (-3,226) rectangle (185,260);
  \node[anchor=base west, inner sep=0pt, text=bkOn, font=\sffamily\fontsize{9.5}{12}\selectfont] at (16,246) {\bkSeries};
  \node[anchor=base west, inner sep=0pt, text=bkOn, font=\sffamily\bfseries\fontsize{19}{24}\selectfont] at (16,233)
    {\kgaBookTitle\hspace{0.5em}{\fontsize{12}{14}\selectfont\kgaBookSub}};
  \node[anchor=base west, inner sep=0pt, text=bkDeep, font=\sffamily\bfseries\fontsize{13}{16}\selectfont] at (16,210) {#1};
  \draw[bkVol, line width=0.8pt] (16,206) -- (166,206);
  #2
  \node[anchor=base west, inner sep=0pt, text=bkDeep, font=\sffamily\bfseries\fontsize{13}{16}\selectfont] at (16,140) {\R{全}{ぜん}9\R{巻}{かん}のならび};
  \draw[bkVol, line width=0.8pt] (16,136) -- (98,136);
  \foreach \i in {1,...,9} {
    \pgfmathsetmacro\bk@ly{127-(\i-1)*8.6}
    \fill[bkVol\i, rounded corners=0.8mm] (16,\bk@ly-2.8) rectangle ++(5.6,5.6);
    \node[text=bkBandOn\i, font=\sffamily\bfseries\fontsize{8}{8}\selectfont] at (18.8,\bk@ly) {\i};
    \node[anchor=west, inner sep=0pt, text=bkInk, font=\sffamily\fontsize{9.5}{11}\selectfont] at (24.5,\bk@ly-0.3)
      {\ltjsetruby{stretch=001}\bkTitle{\i}};
  }
  \node[rectangle callout, callout absolute pointer={(136.5,91)}, anchor=south, draw=bkVol, line width=0.8pt, fill=white,
    rounded corners=3mm, inner xsep=3.5mm, inner ysep=2.5mm, text width=50mm, align=flush center, text=bkDeep,
    font=\sffamily\bfseries\fontsize{10.5}{16}\selectfont] at (136.5,97) {#3};
  \fill[bkVol!22, rounded corners=2mm] (107,50) rectangle (166,54);
  \bkChar{sora}{(117,53.5)}{31mm}\bkChar{mina}{(136.5,53.5)}{31mm}\bkChar{delta}{(156,53.5)}{19mm}
  \draw[bkGray, line width=0.6pt, dash pattern=on 2pt off 1.5pt, fill=white] (114,12) rectangle (166,38);
  \node[text=bkGray, align=flush center, font=\sffamily\fontsize{7}{9}\selectfont] at (140,25) {ISBN・バーコード\\（入稿時に配置）};
  \node[anchor=base west, inner sep=0pt, text=bkGray, font=\sffamily\fontsize{7.5}{10}\selectfont] at (16,25) {\bkSeries　\kgaBookTitle};
  \node[anchor=base west, inner sep=0pt, text=bkGray, font=\sffamily\fontsize{7.5}{10}\selectfont] at (16,16) {\bkEditor\hspace{1pt}\R{編}{へん}};
}
% 表4の中身の1行（札つき）  \kgaBackItem{y}{札の文字}{文}
\newcommand\kgaBackItem[3]{%
  \node[anchor=west, fill=bkVol, text=bkOn, rounded corners=1.2mm, inner xsep=1.6mm, inner ysep=1mm, minimum width=7mm,
    font=\sffamily\bfseries\fontsize{8.5}{10}\selectfont] (kgabi) at (16,#1) {#2};
  \node[anchor=west, inner sep=0pt, text=bkInk, font=\sffamily\fontsize{10}{13}\selectfont] at ([xshift=2.5mm]kgabi.east) {#3};}

%% --------------------------------------------------------------------
%% とびら（p.1）  \kgaTitlePage{kids|adult}
%% --------------------------------------------------------------------
\newcommand\kgaTitleArt[1]{%
  \fill[bkVol] (-3,246) rectangle (185,260);
  \node[anchor=base, text=bkDeep, font=\sffamily\bfseries\fontsize{12}{15}\selectfont] at (91,206) {\bkSeries};
  \kgaBadge{#1}{91}{188}{0.85}
  \node[anchor=base, text=bkDeep, font=\sffamily\bfseries\fontsize{32}{40}\selectfont] at (91,158) {\kgaBookTitle};
  \node[anchor=base, text=bkDeep, font=\sffamily\bfseries\fontsize{15}{18}\selectfont] at (91,145) {\kgaBookSub};
  \node[anchor=base, text=bkGray, font=\fontsize{11}{15}\selectfont] at (91,130) {\kgaCatch};
  \kgaShelf{52.5}{78}{0.5}{bkVol}
  \node[anchor=base, text=bkGray, font=\sffamily\fontsize{9}{12}\selectfont] at (91,34) {\bkEditor\hspace{1pt}\R{編}{へん}};
  \fill[bkVol] (-3,-3) rectangle (185,4);
}
\newcommand\kgaTitlePage[1]{%
  \clearpage\setcounter{page}{1}%
  \begingroup\global\bk@decofalse\bkfolio{false}%
  \AddToHookNext{shipout/background}{\put(0,-\paperheight){%
    \begin{tikzpicture}[overlay, x=1mm, y=1mm]\kgaTitleArt{#1}\end{tikzpicture}}}%
  \null\clearpage\endgroup\global\bk@decotrue\bkfolio{true}}

%% --------------------------------------------------------------------
%% もくじ
%% --------------------------------------------------------------------
% 行  \kgaTocRow{色}{番号}{題名}{ひとこと（空なら1行）}{ラベル}
\newcommand\kgaTocRow[5]{%
  \par\noindent
  \parbox[t]{12mm}{\leavevmode\tikz[baseline=(p.base)]\node[fill=#1, text=white, rounded corners=1.3mm,
      minimum width=9mm, inner xsep=1mm, inner ysep=1.1mm, font=\sffamily\bfseries\fontsize{9.5}{10}\selectfont] (p) {#2};}%
  \parbox[t]{\dimexpr\linewidth-12mm-16mm\relax}{\setlength\parskip{0pt}%
    \leavevmode{\sffamily\bfseries\fontsize{11}{17}\selectfont\color{bkInk}#3}\if\relax\detokenize{#4}\relax\else\par\nobreak\vspace{-0.6mm}
    {\rmfamily\fontsize{9.5}{15}\selectfont\color{bkInk}#4}\fi}%
  \parbox[t]{16mm}{\raggedleft\leavevmode{\sffamily\bfseries\fontsize{11}{17}\selectfont\color{bkInk}\bkpageno{\kgaPage{#5}}}}%
  \par\vspace{1.2mm}{\color{bkLine}\hrule height 0.35pt}\vspace{1.2mm}}
% 見出し  \kgaTocHead
\newcommand\kgaTocHead{%
  \noindent\begin{tikzpicture}
    \node[anchor=base west, inner sep=0pt, text=bkDeep, font=\sffamily\bfseries\fontsize{30}{36}\selectfont] (h) at (0,0) {もくじ};
    \node[anchor=base west, inner sep=0pt, text=bkDeep, font=\sffamily\bfseries\fontsize{12}{15}\selectfont] at ([xshift=6mm]h.base east) {\kgaBookTitle\hspace{0.6em}\kgaBookSub};
    \draw[bkVol, line width=1.6pt, overlay] (0,-3.5mm) -- (\linewidth,-3.5mm);
  \end{tikzpicture}\par\vspace{6mm}}
% \kgaTocPages{1ページめ}{2ページめ（空ならもくじは1ページ）}
\newcommand\kgaTocPages[2]{%
  \clearpage\setcounter{page}{2}%
  \global\bk@decofalse\bkfolio{true}\bksetrunning{}{}{}%
  \kgaTocHead #1\par\clearpage
  \if\relax\detokenize{#2}\relax\else #2\par\clearpage\fi
  \global\bk@decotrue}

%% --------------------------------------------------------------------
%% 色の面の上に問いを置いたページ（子ども用の「はじめに」、p.3）。飾り・ノンブルなし
%%   \kgaArtPage{大きな字}{ラベル}{題名}{問い y/文,…}{絵}{人物}{ふきだしの文}
%% --------------------------------------------------------------------
\newcommand\kgaArtPage[7]{%
  \clearpage
  \begingroup\global\bk@decofalse\bkfolio{false}%
  \AddToHookNext{shipout/background}{\put(0,-\paperheight){%
  \begin{tikzpicture}[overlay, x=1mm, y=1mm]
    \fill[bkVol] (-3,104) rectangle (185,260);
    \begin{scope}\clip (-3,104) rectangle (185,260);
      \node[text=white, opacity=0.12, font=\sffamily\bfseries\fontsize{200}{200}\selectfont] at (140,190) {#1};
    \end{scope}
    \node[anchor=base west, text=white, font=\sffamily\bfseries\fontsize{16}{20}\selectfont] at (22,230) {#2};
    \node[anchor=north west, text=white, font=\sffamily\bfseries\fontsize{30}{40}\selectfont, text width=140mm, align=left, inner sep=0pt] at (21,222) {#3};
    \foreach \y/\t in {#4} {
      \node[anchor=west, inner sep=0pt] at (22,\y) {\bkicon[1.6em]{hatena}{bkYellow}};
      \node[anchor=west, text=white, font=\sffamily\bfseries\fontsize{12.5}{16}\selectfont] at (31,\y) {\t};
    }
    #5
    \bkChar{#6}{(28,62)}{30mm}
    \node[rectangle callout, callout absolute pointer={(43,80)}, anchor=west, draw=bkVol, line width=0.8pt, fill=white,
      rounded corners=3mm, inner xsep=4mm, inner ysep=3mm,
      text width=94mm, align=left, font=\fontsize{10.5}{18}\selectfont, text=bkInk] at (49,80) {#7};
  \end{tikzpicture}}}%
  \null\clearpage\endgroup\global\bk@decotrue\bkfolio{true}}

%% --------------------------------------------------------------------
%% 子ども用の本文の節の見出し（偶数ページから）  \kgasec{番号}{題名}{ひとこと}{しおり用の文字}
%%   ラベル kga-sec番号 を置く（もくじが拾う）。柱：偶数ページ＝\kgaBandLabel、奇数ページ＝番号と題名
%% --------------------------------------------------------------------
\newcommand\kgasec[4]{%
  \bkevenpage
  \global\bk@decotrue
  \bksetrunning{}{\kgaBandLabel}{#1\hspace{0.6\zw}#2}%
  \pdfbookmark[0]{#1　#4}{kgasec#1}\typeout{BKPAGE kga #1 \thepage}%
  \label{kga-sec#1}%
  \noindent\begin{tikzpicture}
    \useasboundingbox (0,0) rectangle (\linewidth,3.1);
    \begin{scope}\clip[rounded corners=2.5mm] (0,0) rectangle (\linewidth,3.1);
      \fill[bkTheme] (0,0) rectangle (\linewidth,3.1);
      \fill[white, opacity=0.13] (\linewidth-2.2,3.1) circle (2.0);
      \fill[white, opacity=0.10] (\linewidth-0.6,0) circle (1.3);
    \end{scope}
    \node[anchor=west, text=bkOn, font=\sffamily\small] at (0.5,2.6) {\kgaBandLabel};
    \node[anchor=west, circle, fill=white, text=\bkTT, minimum size=1.5cm, font=\sffamily\bfseries\Huge, inner sep=0pt] (n) at (0.45,1.25) {#1};
    \node[anchor=west, text=bkOn, font=\sffamily\bfseries\LARGE, text width=\linewidth-3.6cm, align=left, inner sep=0pt] at ([xshift=0.35cm]n.east) {#2};
  \end{tikzpicture}\par\vspace{2.5mm}
  \noindent{\sffamily\fontsize{10.5}{19}\selectfont\color{bkInk}#3}\par
  \vspace{1.5mm}{\color{bkLine}\hrule height 0.5pt}\vspace{3mm}}

%% --------------------------------------------------------------------
%% 奥付  \kgaColophon{書名の行}{副題の行}{版の行}{いちばん下の注}
%% --------------------------------------------------------------------
\newcommand\kgaColophon[4]{%
  \bkpadfour
  \global\bk@decofalse\bkfolio{false}%
  \null\vfill
  \noindent\begin{minipage}{\linewidth}\small
  {\sffamily\bfseries\Large\color{bkVol0}#1}\par\vspace{1mm}
  {\sffamily\bfseries\large #2}\par\vspace{3mm}
  {\color{bkVol0}\hrule height 1pt}\vspace{3mm}
  \begin{tabular}{@{}l@{\hspace{6mm}}l@{}}
  \R{編者}{へん|しゃ} & 「\bkSeries」\R{編集委員会}{へん|しゅう|い|いん|かい} \\
   & \R{委員長}{い|いん|ちょう}　\R{小学校長}{しょう|がっ|こう|ちょう}（まとめ\R{役}{やく}） \\
   & \R{委員}{い|いん}　\R{数学}{すう|がく}・\R{電気工学}{でん|き|こう|がく}・\R{機械工学}{き|かい|こう|がく}・\R{教育学}{きょう|いく|がく}・\R{量子力学}{りょう|し|りき|がく}・\R{建築工学}{けん|ちく|こう|がく}・\\
   & \phantom{委員　}\R{半導体工学}{はん|どう|たい|こう|がく}の\R{大学}{だい|がく}\R{教授}{きょう|じゅ}、\R{幼稚園}{よう|ち|えん}\R{教諭}{きょう|ゆ}、\R{塾}{じゅく}\R{講師}{こう|し}、\\
   & \phantom{委員　}\R{特別支援学校}{とく|べつ|し|えん|がっ|こう}\R{教諭}{きょう|ゆ}、\R{心理士}{しん|り|し}、\R{科学}{か|がく}\R{解説者}{かい|せつ|しゃ}、\R{小学校}{しょう|がっ|こう}\R{教諭}{きょう|ゆ} \\
  デザイン & ユニバーサルデザイン\R{担当}{たん|とう}\R{委員}{い|いん} \\
  \R{挿絵}{さし|え} & イラスト\R{担当}{たん|とう}\R{委員}{い|いん} \\
  \R{版}{はん} & #3 \\
  \end{tabular}\par\vspace{3mm}
  {\color{bkVol0}\hrule height 1pt}\vspace{2mm}
  {\footnotesize #4}
  \end{minipage}
  \clearpage}
\makeatother


\begin{document}
\makeatletter
%% ===================== 前付け =====================
\bkartpage{表1}{\kgaFrontArt{kids}}
\bkInsideFrontPage
\kgaTitlePage{kids}
\kgaTocPages{%
  \noindent{\sffamily\fontsize{10}{17}\selectfont\color{bkGray}この\R{本}{ほん}は、\R{物理}{ぶつ|り}を\R{理解}{り|かい}するための\R{数学}{すう|がく}（\R{全}{ぜん}9\R{巻}{かん}）を\R{読}{よ}み\R{始}{はじ}める\R{前}{まえ}に\R{読}{よ}む\R{本}{ほん}です。ひとりで\R{読}{よ}んでも、おうちの\R{人}{ひと}や\R{先生}{せん|せい}といっしょに\R{読}{よ}んでもかまいません。}\par\vspace{4mm}
  \kgaTocRow{bkVol0}{1}{この\R{本}{ほん}は、なんの\R{本}{ほん}？}{「リンゴは\R{何個}{なん|こ}？」から、ロボット・\R{電気}{でん|き}・\R{音}{おと}まで}{kga-sec1}
  \kgaTocRow{bkVol0}{2}{9\R{冊}{さつ}の\R{本}{ほん}のたな}{1\R{冊}{さつ}ずつ、\R{書名}{しょ|めい}・ページ\R{数}{すう}・★・\R{中身}{なか|み}を\R{見}{み}てみよう}{kga-sec2}
  \kgaTocRow{bkVol0}{3}{7\R{本}{ほん}の\R{路線}{ろ|せん}}{ひき\R{算}{ざん}から\R{微分}{び|ぶん}へ。\R{九九}{く|く}の\R{表}{ひょう}から\R{行列}{ぎょう|れつ}へ。\R{考}{かんが}えは\R{線}{せん}でつながっている}{kga-sec3}
  \kgaTocRow{bkVol0}{4}{1つの\R{章}{しょう}の\R{歩}{ある}き\R{方}{かた}}{はてな？から、ふりかえりまで。ページの\R{部品}{ぶ|ひん}を\R{番号}{ばん|ごう}つきで\R{案内}{あん|ない}}{kga-sec4}
  \kgaTocRow{bkVol0}{5}{★と、3つの\R{段}{だん}}{\R{山}{やま}の\R{高}{たか}さと、「\R{目}{め}→\R{手}{て}→\R{式}{しき}」の3\R{段}{だん}。「ここまでで\R{十分}{じゅう|ぶん}」の\R{旗}{はた}}{kga-sec5}
  \kgaTocRow{bkVol0}{6}{\R{案内役}{あん|ない|やく}のソラ・ミナ・デルタ}{いっしょに\R{歩}{ある}く3\R{人}{にん}を\R{紹介}{しょう|かい}します}{kga-sec6}
  \kgaTocRow{bkVol0}{7}{どこから\R{読}{よ}む？}{はい・いいえで\R{進}{すす}むと、はじめる\R{巻}{かん}が\R{見}{み}つかる\R{流}{なが}れ\R{図}{ず}}{kga-sec7}
  \kgaTocRow{bkVol0}{8}{\R{勉強}{べん|きょう}のしかた}{\R{電卓}{でん|たく}・\R{表計算}{ひょう|けい|さん}・プログラム、こたえの\R{見}{み}\R{方}{かた}、こまったとき}{kga-sec8}
  \kgaTocRow{bkVol0}{9}{\R{巻末}{かん|まつ}の\R{使}{つか}い\R{方}{かた}と、その\R{先}{さき}へ}{\R{記号}{き|ごう}カード・チェック\R{表}{ひょう}・さくいん、そして\R{物理}{ぶつ|り}シリーズへ}{kga-sec9}
  \vfill
  \noindent\begin{tikzpicture}[x=1mm, y=1mm]
    \fill[bkVol0!7, rounded corners=3mm] (0,0) rectangle (140,30);
    \bkChar{delta}{(14,3)}{22mm}
    \node[rectangle callout, callout absolute pointer={(25,17)}, anchor=west, draw=bkVol0, line width=0.8pt, fill=white, rounded corners=3mm,
      inner xsep=3.5mm, inner ysep=2.5mm, text width=96mm, align=left, font=\small, text=bkInk] at (31,15)
      {どの\R{巻}{かん}から\R{始}{はじ}めるか\R{迷}{まよ}ったら、まず7「どこから\R{読}{よ}む？」（\bkpageno{\kgaPage{kga-sec7}}）を\R{開}{ひら}こう。};
  \end{tikzpicture}
}{}
\kgaArtPage{？}{はじめに}{このシリーズの\R{歩}{ある}き\R{方}{かた}}%
  {150/{このシリーズでは、\R{何}{なに}を\R{学}{まな}ぶの？},
   138/{9\R{冊}{さつ}の\R{本}{ほん}は、どうつながっているの？},
   126/{わたしは、どの\R{巻}{かん}から\R{読}{よ}めばいいの？}}%
  {\kgaShelf{112}{108}{0.40}{white}}%
  {delta}%
  {この\R{本}{ほん}は、このシリーズで\R{学}{まな}ぶみんなのための\R{本}{ほん}だよ。9\R{冊}{さつ}の\R{本}{ほん}の\R{中身}{なか|み}と、\R{読}{よ}み\R{方}{かた}と、どこから\R{始}{はじ}めるかを、\R{図}{ず}でたどっていこう。}
\makeatother

%% ===================== 本文 =====================
\IfFileExists{guide-p1.tex}{%% 子ども用ガイドブック『はじめに読む本』本文「このシリーズの歩き方」（第16回編集委員会、組 G-A）。親ファイルは guide-kids.tex
%% 読む本人向け。ルビつき（第2・3巻と同じ \R の付け方。分かち書きはしない）。学年は書かない（決定5）。
%% 数・事実は tools/guide-series.md（RG-2）による。節の見出しは親ファイルの \kgasec（ラベル kga-sec1〜9）。
\makeatletter
%% ===== この本文だけで使う部品（頭文字 kga） =====
% 箇条書きの1項目  \kgaitem{文}
\newcommand\kgaitem[1]{\par\noindent\hangindent=1\zw\hangafter=1 ・#1\par}
% 章番号の札（その章の巻の色）  \kgapill{1-5}
\newcommand\kgapill[1]{\bk@pill{#1}}
% 巻の色の小さな札  \kgavol{巻}
\newcommand\kgavol[1]{\tikz[baseline=(v.base)]\node[fill=bkVol#1, text=bkBandOn#1, rounded corners=0.8mm, inner xsep=1.3mm, inner ysep=0.6mm,
  font=\sffamily\bfseries\footnotesize] (v) {#1};}
% 9冊の本のカード  \kgaVolCard{巻}{ページ}{章の数}{キャッチ}{部の札}{おもな中身}{そのほか}
\newcommand\kgaVolCard[7]{%
  \par\addvspace{1.6mm}\noindent
  \begin{tikzpicture}[x=1mm, y=1mm]
    \node[anchor=north west, text width=101mm, inner xsep=3.5mm, inner ysep=2.2mm, align=left, font=\small, text=bkInk] (c) at (30,0)
      {{\sffamily\bfseries\Large\bkTitle{#1}}\setbox0\hbox{\bkSub{#1}}\ifdim\wd0>0pt\hspace{0.8em}{\sffamily\bfseries\color{bkGray}\bkSub{#1}}\fi\par
       \vspace{0.5mm}{\color{bkGray}#4}\par\vspace{1mm}
       {\sffamily\footnotesize #5}\par\vspace{1mm}
       {\sffamily\bfseries\footnotesize\color{bkInk}おもな\R{中身}{なか|み}}\enspace{\footnotesize #6}\if\relax\detokenize{#7}\relax\else\par\vspace{0.8mm}{\footnotesize\color{bkInk}#7}\fi};
    \draw[bkVol#1, line width=1pt, rounded corners=2.5mm] (0,0) rectangle (c.south east);
    \begin{scope}\clip[rounded corners=2.5mm] (0,0) rectangle (c.south east);
      \fill[bkVol#1] (0,0) rectangle (c.south west);\end{scope}
    \node[text=bkBandOn#1, font=\sffamily\bfseries\footnotesize] at (15,-5) {\R{第}{だい}\hspace{0.15em}#1\hspace{0.15em}\R{巻}{かん}};
    \node[text=bkBandOn#1, font=\sffamily\bfseries\fontsize{34}{34}\selectfont] at (15,-17.5) {#1};
    \node[anchor=north, fill=white, rounded corners=1mm, inner xsep=1.2mm, inner ysep=0.8mm, text=bkInk, font=\sffamily\bfseries\footnotesize] at (15,-27) {#2ページ};
    \node[anchor=north, text=bkBandOn#1, font=\sffamily\bfseries\footnotesize] at (15,-33.5) {#3\R{章}{しょう}};
    \node[anchor=north, fill=white, rounded corners=1mm, inner xsep=1mm, inner ysep=0.8mm] at (15,-40) {\bkstarmeter{\bkSmin{#1}}{\bkSmax{#1}}};
  \end{tikzpicture}\par}
% 部の札  \kgapart{第1部}{名前}
\newcommand\kgapart[2]{\mbox{\tikz[baseline=(p.base)]\node[draw=bkLine, line width=0.6pt, rounded corners=1mm, inner xsep=1.2mm, inner ysep=0.5mm, font=\sffamily\footnotesize, text=bkInk] (p) {{\bfseries #1}\hspace{0.3em}#2};}\hspace{0.3em}\allowbreak}
% 見本のページの番号  \kganum{(x,y)}{番号}
\newcommand\kganum[2]{\node[circle, fill=bkInk, text=white, inner sep=0pt, minimum size=4.6mm, font=\sffamily\bfseries\footnotesize] at #1 {#2};}
% 流れ図の部品
\tikzset{
  kga q/.style={diamond, aspect=3.2, draw=bkTheme, line width=1pt, fill=white, align=flush center, font=\sffamily\fontsize{9.5}{13}\selectfont, text=bkInk, inner sep=0.3mm, text width=42mm},
  kga q1/.style={diamond, aspect=6, draw=bkTheme, line width=1pt, fill=white, align=flush center, font=\sffamily\fontsize{9.5}{13}\selectfont, text=bkInk, inner sep=0.2mm, text width=62mm},
  kga qv/.style={diamond, aspect=4.5, draw=bkTheme, line width=1pt, fill=white, align=flush center, font=\sffamily\fontsize{9.5}{13}\selectfont, text=bkInk, inner sep=0.3mm, text width=50mm},
  kga go/.style={draw=bkVol#1, line width=1.2pt, fill=bkVol#1, text=bkBandOn#1, rounded corners=2mm, align=flush center, font=\sffamily\bfseries\small, inner xsep=2mm, inner ysep=1.6mm, minimum width=23mm},
  kga box/.style={draw=bkGray, line width=0.9pt, fill=white, rounded corners=2mm, align=flush center, font=\sffamily\small, text=bkInk, inner xsep=2.5mm, inner ysep=2mm},
  kga arr/.style={-{Stealth[length=2.4mm]}, line width=1.1pt, bkGray, shorten <=1.5mm, shorten >=1.5mm},
  kga yn/.style={font=\sffamily\bfseries\footnotesize, text=bkInk, inner sep=0.8mm, fill=white},
  kga lab/.style={font=\sffamily\footnotesize, text=bkInk, inner sep=0.5mm},
}
% 会話の1行（2-0 と同じ形）  \kgatalk{y}{sora|mina|delta}{名前}{せりふ}
\newcommand\kgatalk[4]{%
  \bkChar{#2}{(7,#1-6.5)}{12mm}
  \node[rectangle callout, callout absolute pointer={(15.5,#1)}, anchor=west, draw=bkTheme, line width=0.8pt,
    fill=white, rounded corners=2.5mm, inner xsep=3mm, inner ysep=1.8mm, text width=103mm, align=left,
    font=\small] at (21,#1) {#4};
  \node[anchor=north, font=\sffamily\footnotesize, text=bkGray, inner sep=0pt] at (7,#1-7.5) {#3};}
\makeatother

%% =====================================================================
%% 1　この本は、なんの本？（2ページ）
%% =====================================================================
\kgasec{1}{この\R{本}{ほん}は、なんの\R{本}{ほん}？}
  {「リンゴは\R{何個}{なん|こ}？」から\R{始}{はじ}めて、ロボット・\R{電気}{でん|き}・\R{音}{おと}・\R{機械}{き|かい}の\R{動}{うご}きを\R{読}{よ}みとく\R{数学}{すう|がく}まで、9\R{冊}{さつ}でつながっています。}
  {この本は、なんの本？}

このシリーズ『\bkSeries』は、\R{全}{ぜん}9\R{巻}{かん}の\R{算数}{さん|すう}・\R{数学}{すう|がく}の\R{本}{ほん}です。\R{第}{だい}1\R{巻}{かん}は、「リンゴは\R{何個}{なん|こ}？」のように\R{数}{かぞ}えるところから\R{始}{はじ}まります。\R{第}{だい}9\R{巻}{かん}では、\R{音}{おと}をいくつもの\R{波}{なみ}に\R{分}{わ}けたり、\R{機械}{き|かい}や\R{電気}{でん|き}の\R{回路}{かい|ろ}がこの\R{先}{さき}どう\R{動}{うご}くかを\R{予想}{よ|そう}したりする\R{数学}{すう|がく}まで\R{進}{すす}みます。この\R{本}{ほん}（ガイドブック）は、その9\R{冊}{さつ}を\R{読}{よ}み\R{始}{はじ}める\R{前}{まえ}に\R{読}{よ}む\R{本}{ほん}です。

\begin{bkfigure}
\begin{tikzpicture}[x=1mm, y=1mm]
  \foreach \v in {1,...,9} {
    \pgfmathsetmacro\kgax{4+(\v-1)*14.8}\pgfmathsetmacro\kgah{8+(\v-1)*4.4}
    \fill[bkVol\v, rounded corners=0.8mm] (\kgax,0) rectangle ++(13.6,\kgah);
    \node[text=bkBandOn\v, font=\sffamily\bfseries\large] at (\kgax+6.8,\kgah-4.5) {\v};}
  \draw[bkInk!60, line width=0.8pt] (0,0) -- (140,0);
  \bkapple{(10.8,13)}{3.2}{bkAppleRed}
  \tikzset{kgacl/.style={draw=bkGray, line width=0.7pt, fill=white, rounded corners=1.5mm, align=left, font=\sffamily\footnotesize, text=bkInk, inner xsep=1.6mm, inner ysep=1.2mm, anchor=south west}}
  \node[kgacl] at (0,19) {\kgapill{1-0}\enspace かずを つかうと、なにが いいの？};
  \node[kgacl] at (0,27.5) {\kgapill{2-0}\enspace「だんだん\R{速}{はや}く」を\R{数}{かず}で\R{言}{い}うと？};
  \node[kgacl] at (0,36) {\kgapill{4-0}\enspace ロボットの\R{手}{て}は、どこにある？};
  \node[kgacl] at (0,44.5) {\kgapill{6-0}\enspace「\R{次}{つぎ}の\R{瞬間}{しゅん|かん}を\R{決}{き}めるルール」で\R{未来}{み|らい}を\R{予想}{よ|そう}する};
  \node[kgacl] at (0,53) {\kgapill{9-0}\enspace\R{音}{おと}を\R{分}{わ}ける、\R{変化}{へん|か}を\R{予想}{よ|そう}する};
\end{tikzpicture}
\bkfigcaption{9\R{冊}{さつ}は1\R{本}{ぽん}の\R{階段}{かい|だん}。\R{左上}{ひだり|うえ}は\R{序章}{じょ|しょう}（\R{巻}{かん}のはじめの\R{章}{しょう}）の\R{題}{だい}の\R{例}{れい}で、\R{札}{ふだ}の\R{色}{いろ}は\R{巻}{かん}の\R{色}{いろ}です。}
\end{bkfigure}

\bksubsection{1}{なぜ「\R{物理}{ぶつ|り}を\R{理解}{り|かい}するための\R{数学}{すう|がく}」なの？}

\R{物理}{ぶつ|り}は、ものの\R{動}{うご}き・\R{熱}{ねつ}・\R{電気}{でん|き}・\R{光}{ひかり}や\R{音}{おと}のきまりを\R{調}{しら}べる\R{学問}{がく|もん}、\R{工学}{こう|がく}は、そのきまりを\R{使}{つか}って\R{機械}{き|かい}・\R{電気}{でん|き}の\R{回路}{かい|ろ}・ロボットなどを\R{作}{つく}る\R{学問}{がく|もん}です。このシリーズでは、\R{入試}{にゅう|し}の\R{問題}{もん|だい}を\R{解}{と}くための\R{数学}{すう|がく}よりも、\R{物理}{ぶつ|り}や\R{工学}{こう|がく}のできごとを\R{説明}{せつ|めい}するための\R{数学}{すう|がく}を\R{学}{まな}びます。\R{最後}{さい|ご}まで\R{学}{まな}ぶと、\R{大学}{だい|がく}の\R{工学}{こう|がく}の\R{専門書}{せん|もん|しょ}の\R{式}{しき}が、すんなり\R{頭}{あたま}に\R{入}{はい}るようになることを\R{目指}{め|ざ}しています。

\bksubsection{2}{1つのできごとを、\R{巻}{かん}を\R{進}{すす}むごとにくわしく}

たとえば「\R{冷}{さ}めるお\R{湯}{ゆ}」は、いくつもの\R{巻}{かん}に\R{出}{で}てきます。

\begin{bkfigure}
\begin{tikzpicture}[x=1mm, y=1mm]
  \foreach \c/\t [count=\i from 0] in {2-18/{\R{表計算}{ひょう|けい|さん}で、1\R{分}{ぷん}ずつ\R{予想}{よ|そう}する},
      3-11/{「\R{同}{おな}じ\R{割合}{わり|あい}でへる」\R{指数関数}{し|すう|かん|すう}},
      5-14/{\R{数}{かず} $e$ と\R{時定数}{じ|てい|すう}},
      6-4/{\R{微分方程式}{び|ぶん|ほう|てい|しき}を\R{立}{た}てて\R{解}{と}く},
      9-0/{\R{変換表}{へん|かん|ひょう}を\R{引}{ひ}いて\R{予想}{よ|そう}する}} {
    \node[draw=bkTheme, line width=0.9pt, fill=white, rounded corners=2mm, minimum width=25.5mm, minimum height=21mm, anchor=north west] (k\i) at (\i*28.6,0) {};
    \node[anchor=north] at ([yshift=-1.8mm]k\i.north) {\kgapill{\c}};
    \node[anchor=north, text width=23mm, align=flush center, font=\sffamily\footnotesize, text=bkInk, inner sep=0pt] at ([yshift=-7.6mm]k\i.north) {\t};}
  \foreach \a/\b in {k0/k1,k1/k2,k2/k3,k3/k4} \draw[-{Stealth[length=2mm]}, line width=1pt, bkGray] ([xshift=0.4mm]\a.east) -- ([xshift=-0.4mm]\b.west);
\end{tikzpicture}
\bkfigcaption{「\R{冷}{さ}めるお\R{湯}{ゆ}」が\R{出}{で}てくる\R{章}{しょう}。\R{札}{ふだ}の\R{色}{いろ}は\R{巻}{かん}の\R{色}{いろ}です。}
\end{bkfigure}

\noindent\begin{minipage}[c]{0.52\linewidth}\centering
\begin{tikzpicture}[x=1mm, y=1mm]
  % 横：0〜10分（6 mm/分）、たて：50〜80 ℃（1.25 mm/℃）
  \foreach \t in {0,2,...,10} {\draw[bkLine, line width=0.4pt] (\t*6,0) -- (\t*6,30);
    \draw[bkInk, line width=0.6pt] (\t*6,-8) -- (\t*6,-9);
    \node[anchor=north, font=\sffamily\footnotesize, text=bkInk, inner sep=1pt] at (\t*6,-9.2) {\t};}
  \foreach \T in {50,60,70,80} {\draw[bkLine, line width=0.4pt] (0,{(\T-50)*1.0}) -- (60,{(\T-50)*1.0});
    \node[anchor=east, font=\sffamily\footnotesize, text=bkInk, inner sep=1pt] at (-0.8,{(\T-50)*1.0}) {\T};}
  \draw[bkInk, line width=0.8pt] (0,-8) -- (0,33) (0,-8) -- (63,-8);
  \node[anchor=east, font=\sffamily\footnotesize, text=bkInk, inner sep=1pt] at (-0.8,-8) {0};
  % 0 から 50 までを省いた印（波線）
  \fill[white] (-2,-5.2) rectangle (2,-2.8);
  \draw[bkInk, line width=0.7pt] (-2,-4.6) .. controls (-1,-3.6) and (0,-5.6) .. (1,-4.6) -- (2,-4.1);
  \draw[bkInk, line width=0.7pt] (-2,-3.4) .. controls (-1,-2.4) and (0,-4.4) .. (1,-3.4) -- (2,-2.9);
  \node[anchor=south west, font=\sffamily\footnotesize, text=bkInk, inner sep=0pt] at (-6,34) {\R{温度}{おん|ど}（℃）};
  \node[anchor=north east, font=\sffamily\footnotesize, text=bkInk, inner sep=0pt] at (63,-14) {\R{時間}{じ|かん}（\R{分}{ふん}）};
  \draw[bkBlue, line width=1.5pt, domain=0:10, samples=60, smooth] plot ({\x*6},{(60*exp(-0.05*\x)-30)*1.0});
  \foreach \t/\T in {0/80,1/77,2/74.15,3/71.44,4/68.87,5/66.43,6/64.11,7/61.90,8/59.81,9/57.81,10/55.92}
    {\fill[white] (\t*6,{(\T-50)*1.0}) circle (1.05); \draw[bkOrange!85!black, line width=0.9pt] (\t*6,{(\T-50)*1.0}) circle (1.05);}
\end{tikzpicture}
\end{minipage}\hfill
\begin{minipage}[c]{0.45\linewidth}
\small
\R{部屋}{へ|や}の\R{温度}{おん|ど}が \qty{20}{\degreeCelsius} のとき、\qty{80}{\degreeCelsius} のお\R{湯}{ゆ}の\R{冷}{さ}め\R{方}{かた}です。\par\vspace{1mm}
\tikz[baseline=-0.6ex]{\fill[white] (0,0) circle (1.05mm); \draw[bkOrange!85!black, line width=0.9pt] (0,0) circle (1.05mm);}\enspace\R{表計算}{ひょう|けい|さん}で1\R{分}{ぷん}ずつ\R{進}{すす}めた\R{値}{あたい}（10\R{分}{ぷん}で\R{約}{やく}\qty{55.9}{\degreeCelsius}）\par\vspace{0.5mm}
\tikz[baseline=-0.6ex]\draw[bkBlue, line width=1.5pt] (0,0) -- (5mm,0);\enspace\R{微分方程式}{び|ぶん|ほう|てい|しき}を\R{解}{と}いた\R{式}{しき}の\R{値}{あたい}（10\R{分}{ぷん}で\R{約}{やく}\qty{56.4}{\degreeCelsius}）
\end{minipage}
\par\vspace{1mm}
{\centering\sffamily\small{\color{\bkTT}\bfseries \R{図}{ず}\refstepcounter{bkfig}\thebkfig}\hspace{0.8\zw}\R{表}{ひょう}で\R{予想}{よ|そう}した\R{点}{てん}と、\R{式}{しき}で\R{求}{もと}めた\R{線}{せん}は、ほとんど\R{重}{かさ}なる\par}\vspace{3mm}

\bksubsection{3}{このシリーズの\R{学}{まな}び\R{方}{かた}}

\begin{bkfigure}
\begin{tikzpicture}[x=1mm, y=1mm]
  \foreach \t/\s [count=\i from 0] in {\R{場面}{ば|めん}/{\R{見}{み}る・ふしぎに\R{思}{おも}う}, \R{表}{ひょう}/{はかる・\R{数}{かぞ}える}, グラフ/{\R{点}{てん}を\R{打}{う}って\R{形}{かたち}を\R{見}{み}る}, \R{式}{しき}/{\R{表}{ひょう}とグラフを\R{短}{みじか}く\R{書}{か}く}} {
    \node[draw=bkTheme, line width=1pt, fill=bkTheme!6, rounded corners=2mm, minimum width=27mm, minimum height=8.5mm, anchor=west, font=\sffamily\bfseries\small, text=\bkTT] (s\i) at (\i*35.5,0) {\t};}
  \foreach \a/\b in {s0/s1,s1/s2,s2/s3} \draw[-{Stealth[length=2.2mm]}, line width=1.1pt, bkGray] ([xshift=0.8mm]\a.east) -- ([xshift=-0.8mm]\b.west);
\end{tikzpicture}
\bkfigcaption{\R{場面}{ば|めん}を\R{見}{み}て、はかって\R{表}{ひょう}にし、グラフにかいてから\R{式}{しき}にする（\R{式}{しき}は\R{最後}{さい|ご}）}
\end{bkfigure}

\kgaitem{\R{数}{かず}には、いつも\R{単位}{たん|い}をつけます。「5」ではなく「\qty{5}{m}」です。}
\kgaitem{\R{新}{あたら}しいことばや\R{記号}{き|ごう}は、\R{手}{て}を\R{動}{うご}かして\R{気}{き}づいたあとで\R{紹介}{しょう|かい}します。}
\kgaitem{くり\R{返}{かえ}しの\R{計算}{けい|さん}は、\R{第}{だい}2\R{巻}{かん}から\R{電卓}{でん|たく}やコンピュータにまかせます。}

%% =====================================================================
%% 2　9冊の本のたな（4ページ）
%% =====================================================================
\kgasec{2}{9\R{冊}{さつ}の\R{本}{ほん}のたな}
  {\R{全}{ぜん}9\R{巻}{かん}、あわせて4836ページ、152の\R{章}{しょう}。1\R{冊}{さつ}ずつ、\R{書名}{しょ|めい}・ページ\R{数}{すう}・★・\R{中身}{なか|み}を\R{見}{み}てみよう。}
  {9冊の本のたな}

このシリーズは、9\R{冊}{さつ}の\R{本}{ほん}でできています。\R{本}{ほん}ごとに\R{表紙}{ひょう|し}の\R{色}{いろ}がちがい、\R{章}{しょう}の\R{番号}{ばん|ごう}の\R{札}{ふだ}も、その\R{巻}{かん}の\R{色}{いろ}でぬってあります（「3-5」は\R{第}{だい}3\R{巻}{かん}の5\R{章}{しょう}）。\R{色}{いろ}だけで\R{見分}{み|わ}けなくてもいいように、\R{札}{ふだ}にはかならず\R{巻}{かん}の\R{番号}{ばん|ごう}が\R{書}{か}いてあります。

\begin{bkfigure}
\begin{tikzpicture}[x=1mm, y=1mm]
  \foreach \v/\p/\c in {1/400/22,2/456/20,3/564/18,4/560/16,5/492/19,6/600/16,7/492/12,8/632/15,9/640/14} {
    \pgfmathsetmacro\kgax{6+(\v-1)*14.6}\pgfmathsetmacro\kgah{\p*0.1}
    \fill[bkVol\v, rounded corners=0.8mm] (\kgax,0) rectangle ++(12.6,\kgah);
    \draw[bkBandOn\v, line width=0.5pt, opacity=0.6] (\kgax+1.4,\kgah-4) -- ++(9.8,0) (\kgax+1.4,4) -- ++(9.8,0);
    \node[text=bkBandOn\v, font=\sffamily\bfseries\large] at (\kgax+6.3,\kgah-10) {\v};
    \node[anchor=south, font=\sffamily\footnotesize, text=bkInk, inner sep=0.8pt] at (\kgax+6.3,\kgah+0.6) {\p};
    \node[anchor=north, font=\sffamily\footnotesize, text=bkInk, inner sep=0.8pt] at (\kgax+6.3,-2.8) {\c\R{章}{しょう}};}
  \fill[bkStem!70, rounded corners=0.6mm] (2,-2.4) rectangle (139,0);
  \node[anchor=south west, font=\sffamily\footnotesize, text=bkGray, inner sep=0pt] at (0,68) {\R{本}{ほん}の\R{上}{うえ}の\R{数}{かず}＝ページ\R{数}{すう}};
\end{tikzpicture}
\bkfigcaption{\R{本}{ほん}の\R{高}{たか}さはページ\R{数}{すう}に\R{合}{あ}わせてかいた。\R{下}{した}は\R{章}{しょう}の\R{数}{かず}（\R{序章}{じょ|しょう}・\R{読}{よ}みもの・\R{終章}{しゅう|しょう}をふくむ）}
\end{bkfigure}

\bksubsection{1}{カードの\R{見}{み}\R{方}{かた}}

\R{次}{つぎ}のページから、1\R{冊}{さつ}ずつカードで\R{紹介}{しょう|かい}します。\R{左}{ひだり}の\R{色}{いろ}の\R{帯}{おび}は\R{表紙}{ひょう|し}の\R{色}{いろ}で、\R{巻}{かん}の\R{番号}{ばん|ごう}・ページ\R{数}{すう}・\R{章}{しょう}の\R{数}{かず}が\R{書}{か}いてあります。いちばん\R{下}{した}の★は、その\R{巻}{かん}の★の\R{範囲}{はん|い}です。\R{右}{みぎ}には、\R{書名}{しょ|めい}と\R{副題}{ふく|だい}、\R{表紙}{ひょう|し}のことば、\R{部}{ぶ}の\R{名前}{な|まえ}、おもな\R{中身}{なか|み}がならんでいます。★の\R{読}{よ}み\R{方}{かた}は5\R{節}{せつ}で\R{説明}{せつ|めい}します。

\clearpage
\kgaVolCard{1}{400}{22}{\bkCatch{1}}
  {\kgapart{\R{第}{だい}1\R{部}{ぶ}}{かずの はじまり}\kgapart{\R{第}{だい}2\R{部}{ぶ}}{はかる}\kgapart{\R{第}{だい}3\R{部}{ぶ}}{かけ\R{算}{ざん}と わり\R{算}{ざん}}\kgapart{\R{第}{だい}4\R{部}{ぶ}}{\R{表}{ひょう}・グラフ・かたち}}
  {\R{数}{かぞ}え\R{方}{かた}と10のまとまり、たし\R{算}{ざん}・ひき\R{算}{ざん}と\R{筆算}{ひっ|さん}、\R{長}{なが}さ・かさ・\R{重}{おも}さ・\R{時刻}{じ|こく}、かけ\R{算}{ざん}と\R{九九}{く|く}、わり\R{算}{ざん}、\R{小数}{しょう|すう}と\R{分数}{ぶん|すう}の\R{入口}{いり|ぐち}、\R{表}{ひょう}とグラフ、\R{形}{かたち}、\R{角}{かく}と\R{面積}{めん|せき}}
  {\R{分}{わ}かち\R{書}{が}き・ふりがなつき。\R{問題}{もん|だい}は\R{本}{ほん}に\R{書}{か}きこみます。かけ\R{算}{ざん}・わり\R{算}{ざん}の\R{筆算}{ひっ|さん}は、\R{学校}{がっ|こう}の\R{教科書}{きょう|か|しょ}で\R{練習}{れん|しゅう}しよう。}
\kgaVolCard{2}{456}{20}{\bkCatch{2}}
  {\kgapart{\R{第}{だい}1\R{部}{ぶ}}{\R{量}{りょう}と\R{数}{かず}}\kgapart{\R{第}{だい}2\R{部}{ぶ}}{\R{大}{おお}きな\R{数}{かず}・\R{小}{ちい}さな\R{数}{かず}・\R{単位}{たん|い}}\kgapart{\R{第}{だい}3\R{部}{ぶ}}{\R{文字}{も|じ}と\R{式}{しき}}\kgapart{\R{第}{だい}4\R{部}{ぶ}}{コンピュータで\R{計算}{けい|さん}}}
  {\R{量}{りょう}＝\R{数}{かず}×\R{単位}{たん|い}、マイナスの\R{数}{かず}、\R{分数}{ぶん|すう}と\R{小数}{しょう|すう}の\R{計算}{けい|さん}、\R{比}{ひ}と\R{割合}{わり|あい}、はかった\R{数}{かず}のずれ、\R{累乗}{るい|じょう}・\R{平方根}{へい|ほう|こん}・\R{円周率}{えん|しゅう|りつ}、SI\R{単位}{たん|い}、\R{文字}{も|じ}と\R{式}{しき}、\R{比例}{ひ|れい}と\R{反比例}{はん|ぴ|れい}、1\R{次方程式}{じ|ほう|てい|しき}、2\R{進数}{しん|すう}、\R{表計算}{ひょう|けい|さん}とプログラム}
  {ふりがなつき。ここから、\R{問題}{もん|だい}はノートに\R{解}{と}きます。\R{電卓}{でん|たく}は 2-1 から\R{使}{つか}います。}
\kgaVolCard{3}{564}{18}{\bkCatch{3}}
  {\kgapart{\R{第}{だい}1\R{部}{ぶ}}{\R{関数}{かん|すう}とグラフ}\kgapart{\R{第}{だい}2\R{部}{ぶ}}{\R{曲}{ま}がったグラフ}\kgapart{\R{第}{だい}3\R{部}{ぶ}}{\R{倍々}{ばい|ばい}とけた}\kgapart{\R{第}{だい}4\R{部}{ぶ}}{データと\R{統計}{とう|けい}}}
  {\R{変数}{へん|すう}と\R{座標}{ざ|ひょう}、\R{関数}{かん|すう}、グラフの\R{傾}{かたむ}き、\R{連立方程式}{れん|りつ|ほう|てい|しき}、\R{展開}{てん|かい}と\R{平方完成}{へい|ほう|かん|せい}、2\R{次関数}{じ|かん|すう}と2\R{次方程式}{じ|ほう|てい|しき}、\R{指数関数}{し|すう|かん|すう}と\R{対数}{たい|すう}、\R{平均}{へい|きん}と\R{標準偏差}{ひょう|じゅん|へん|さ}、\R{相関}{そう|かん}と\R{最小二乗法}{さい|しょう|に|じょう|ほう}、\R{確率}{かく|りつ}と\R{正規分布}{せい|き|ぶん|ぷ}}
  {ふりがなつきは、この\R{巻}{かん}まで。}

\clearpage
\kgaVolCard{4}{560}{16}{\R{長}{なが}さと\R{角度}{かく|ど}で、\R{形}{かたち}と\R{向}{む}きを\R{数}{かず}にしよう}
  {\kgapart{\R{第}{だい}1\R{部}{ぶ}}{\R{図形}{ず|けい}}\kgapart{\R{第}{だい}2\R{部}{ぶ}}{\R{三角比}{さん|かく|ひ}と\R{三角関数}{さん|かく|かん|すう}}\kgapart{\R{第}{だい}3\R{部}{ぶ}}{ベクトル}}
  {\R{定義}{てい|ぎ}・\R{定理}{てい|り}と\R{証明}{しょう|めい}、\R{角}{かく}と\R{円}{えん}・ラジアン、\R{相似}{そう|じ}、\R{三平方}{さん|へい|ほう}の\R{定理}{てい|り}、\R{面積}{めん|せき}と\R{体積}{たい|せき}、\R{三角比}{さん|かく|ひ}、\R{正弦定理}{せい|げん|てい|り}と\R{余弦定理}{よ|げん|てい|り}、\R{三角関数}{さん|かく|かん|すう}と\R{加法定理}{か|ほう|てい|り}、ベクトル・\R{内積}{ない|せき}・\R{外積}{がい|せき}、\R{空間}{くう|かん}の\R{座標}{ざ|ひょう}、いろいろな\R{座標}{ざ|ひょう}}
  {ここからは、ふりがなはありません（\R{使}{つか}う\R{漢字}{かん|じ}は\R{常用漢字}{じょう|よう|かん|じ}まで）。}
\kgaVolCard{5}{492}{19}{\R{変}{か}わる\R{速}{はや}さと、たまった\R{量}{りょう}を\R{計算}{けい|さん}しよう}
  {\kgapart{\R{第}{だい}1\R{部}{ぶ}}{\R{微分}{び|ぶん}}\kgapart{\R{第}{だい}2\R{部}{ぶ}}{\R{積分}{せき|ぶん}}\kgapart{\R{第}{だい}3\R{部}{ぶ}}{\R{特別}{とく|べつ}な\R{関数}{かん|すう}}}
  {\R{平均}{へい|きん}の\R{速}{はや}さと\R{瞬間}{しゅん|かん}の\R{速}{はや}さ、\R{極限}{きょく|げん}と\R{微分}{び|ぶん}、\R{位置}{い|ち}・\R{速度}{そく|ど}・\R{加速度}{か|そく|ど}、\R{最大}{さい|だい}と\R{最小}{さい|しょう}、かけ\R{算}{ざん}での\R{近似}{きん|じ}、\R{細}{ほそ}い\R{長方形}{ちょう|ほう|けい}に\R{分}{わ}けてたす\R{積分}{せき|ぶん}、\R{積分}{せき|ぶん}は\R{微分}{び|ぶん}の\R{逆}{ぎゃく}、\R{仕事}{し|ごと}と\R{流量}{りゅう|りょう}、\R{数}{かず} $e$、sin・cos の\R{微分}{び|ぶん}、テイラー\R{展開}{てん|かい}}
  {}
\kgaVolCard{6}{600}{16}{「\R{次}{つぎ}の\R{瞬間}{しゅん|かん}のルール」から、\R{未来}{み|らい}を\R{計算}{けい|さん}しよう}
  {\kgapart{\R{第}{だい}1\R{部}{ぶ}}{\R{微分方程式}{び|ぶん|ほう|てい|しき}の\R{入口}{いり|ぐち}}\kgapart{\R{第}{だい}2\R{部}{ぶ}}{\R{複素数}{ふく|そ|すう}と\R{双曲線関数}{そう|きょく|せん|かん|すう}}\kgapart{\R{第}{だい}3\R{部}{ぶ}}{\R{式}{しき}で\R{解}{と}く}}
  {\R{微分方程式}{び|ぶん|ほう|てい|しき}とオイラー\R{法}{ほう}（\R{表計算}{ひょう|けい|さん}で\R{解}{と}く）、\R{冷}{さ}めるお\R{湯}{ゆ}と\R{時定数}{じ|てい|すう}、ばねと\R{振}{ふ}り\R{子}{こ}、ゆれながら\R{止}{と}まる\R{動}{うご}き、\R{複素数}{ふく|そ|すう}と\R{複素平面}{ふく|そ|へい|めん}、オイラーの\R{公式}{こう|しき}、\R{双曲線関数}{そう|きょく|せん|かん|すう}、\R{特性方程式}{とく|せい|ほう|てい|しき}、\R{共振}{きょう|しん}、\R{連立微分方程式}{れん|りつ|び|ぶん|ほう|てい|しき}}
  {7\R{本}{ほん}の\R{路線}{ろ|せん}のうち4\R{本}{ほん}が、この\R{巻}{かん}で\R{乗}{の}りかえます（3\R{節}{せつ}）。}

\clearpage
\kgaVolCard{7}{492}{12}{たくさんの\R{数}{かず}を、\R{表}{ひょう}のまま\R{計算}{けい|さん}しよう}
  {\kgapart{\R{第}{だい}1\R{部}{ぶ}}{\R{行列}{ぎょう|れつ}の\R{計算}{けい|さん}}\kgapart{\R{第}{だい}2\R{部}{ぶ}}{\R{回転}{かい|てん}と\R{座標}{ざ|ひょう}の\R{変換}{へん|かん}}\kgapart{\R{第}{だい}3\R{部}{ぶ}}{\R{固有値}{こ|ゆう|ち}と\R{行列}{ぎょう|れつ}の\R{使}{つか}いみち}}
  {\R{行列}{ぎょう|れつ}の\R{計算}{けい|さん}、\R{形}{かたち}を\R{変}{か}える\R{行列}{ぎょう|れつ}（1\R{次変換}{じ|へん|かん}）、\R{行列式}{ぎょう|れつ|しき}、\R{逆行列}{ぎゃく|ぎょう|れつ}で\R{連立方程式}{れん|りつ|ほう|てい|しき}を\R{解}{と}く、\R{回転行列}{かい|てん|ぎょう|れつ}、3\R{次元}{じ|げん}の\R{回転}{かい|てん}と\R{座標}{ざ|ひょう}の\R{変換}{へん|かん}、\R{固有値}{こ|ゆう|ち}と\R{固有}{こ|ゆう}ベクトル、\R{対角化}{たい|かく|か}、\R{行列}{ぎょう|れつ}のくり\R{返}{かえ}し、\R{読}{よ}みもの「\R{行列}{ぎょう|れつ}の\R{使}{つか}いみち」}
  {}
\kgaVolCard{8}{632}{15}{\R{空間}{くう|かん}に\R{広}{ひろ}がる\R{量}{りょう}を、\R{矢印}{や|じるし}で\R{読}{よ}みとこう}
  {\kgapart{\R{第}{だい}1\R{部}{ぶ}}{たくさんの\R{変数}{へん|すう}の\R{微分}{び|ぶん}}\kgapart{\R{第}{だい}2\R{部}{ぶ}}{ベクトル\R{場}{ば}と\R{積分}{せき|ぶん}}\kgapart{\R{第}{だい}3\R{部}{ぶ}}{\R{発散}{はっ|さん}と\R{回転}{かい|てん}}\kgapart{\R{第}{だい}4\R{部}{ぶ}}{\R{偏微分方程式}{へん|び|ぶん|ほう|てい|しき}}}
  {\R{等高線}{とう|こう|せん}と\R{偏微分}{へん|び|ぶん}・\R{全微分}{ぜん|び|ぶん}、\R{勾配}{こう|ばい}、ベクトル\R{場}{ば}、\R{重積分}{じゅう|せき|ぶん}・\R{線積分}{せん|せき|ぶん}・\R{面積分}{めん|せき|ぶん}、\R{発散}{はっ|さん}（div）と\R{回転}{かい|てん}（rot）、ガウスの\R{定理}{てい|り}とストークスの\R{定理}{てい|り}、ラプラシアン、\R{熱}{ねつ}と\R{波}{なみ}の\R{偏微分方程式}{へん|び|ぶん|ほう|てい|しき}、\R{読}{よ}みもの「div と rot で\R{書}{か}かれた\R{式}{しき}を\R{読}{よ}もう」}
  {}
\kgaVolCard{9}{640}{14}{\R{波}{なみ}を\R{成分}{せい|ぶん}に\R{分}{わ}け、\R{微分方程式}{び|ぶん|ほう|てい|しき}をかけ\R{算}{ざん}で\R{解}{と}こう}
  {\kgapart{\R{第}{だい}1\R{部}{ぶ}}{フーリエ\R{解析}{かい|せき}}\kgapart{\R{第}{だい}2\R{部}{ぶ}}{ラプラス\R{変換}{へん|かん}}\kgapart{\R{第}{だい}3\R{部}{ぶ}}{\R{行列}{ぎょう|れつ}と\R{微分方程式}{び|ぶん|ほう|てい|しき}}\kgapart{\R{終章}{しゅう|しょう}}{\R{工学}{こう|がく}の\R{数学}{すう|がく}の\R{地図}{ち|ず}}}
  {フーリエ\R{級数}{きゅう|すう}・フーリエ\R{変換}{へん|かん}・FFT、ラプラス\R{変換}{へん|かん}と\R{変換表}{へん|かん|ひょう}、\R{伝達関数}{でん|たつ|かん|すう}、\R{極}{きょく}と\R{安定性}{あん|てい|せい}、\R{状態方程式}{じょう|たい|ほう|てい|しき}、\R{行列}{ぎょう|れつ}の\R{指数関数}{し|すう|かん|すう}、\R{読}{よ}みもの「フーリエとラプラスの\R{使}{つか}いみち」}
  {\R{終章}{しゅう|しょう}で、7\R{本}{ほん}の\R{路線}{ろ|せん}がどこで\R{合流}{ごう|りゅう}したかをたどります。}

%% =====================================================================
%% 3　7本の路線（2ページ）
%% =====================================================================
\kgasec{3}{7\R{本}{ほん}の\R{路線}{ろ|せん}}
  {ひき\R{算}{ざん}は\R{微分}{び|ぶん}へ、\R{九九}{く|く}の\R{表}{ひょう}は\R{行列}{ぎょう|れつ}へ。\R{考}{かんが}えは、\R{巻}{かん}をこえて\R{線}{せん}でつながっています。}
  {7本の路線}

このシリーズの\R{章}{しょう}は、どれも7\R{本}{ほん}の\R{路線}{ろ|せん}のどれかの「\R{駅}{えき}」になっています。\R{同}{おな}じ\R{路線}{ろ|せん}の\R{駅}{えき}は、1つの\R{考}{かんが}え\R{方}{かた}でつながっています。たとえば\R{変化}{へん|か}の\R{路線}{ろ|せん}では、\R{第}{だい}1\R{巻}{かん}の「ひき\R{算}{ざん}は、かわった\R{量}{りょう}」（1-5）が、\R{第}{だい}3\R{巻}{かん}でグラフの\R{傾}{かたむ}き（3-5）になり、\R{第}{だい}5\R{巻}{かん}で\R{微分}{び|ぶん}（5-1）につながります。

\begin{bkcenter}
\begin{bktable}{colspec={Q[c,m,wd=7mm] Q[l,m,wd=25mm] Q[l,m,wd=31mm] Q[l,m,wd=38mm] Q[l,m,wd=19mm]}, colsep=4pt, rowsep=1.6pt, cells={font=\small}, row{1}={font=\sffamily\bfseries\small}}
 & \R{路線}{ろ|せん} & \R{流}{なが}れ & おもな\R{駅}{えき}（\R{左}{ひだり}がはじまり） & \R{終点}{しゅう|てん} \\
\bkrouteicon{1} & ①\,\R{変化}{へん|か} & {ひき\R{算}{ざん}→\R{傾}{かたむ}き\\→\R{微分}{び|ぶん}} & \kgapill{1-5}\hspace{0.25em}\allowbreak\kgapill{2-3}\hspace{0.25em}\allowbreak\kgapill{3-5}\hspace{0.25em}\allowbreak\kgapill{5-1}\hspace{0.25em}\allowbreak\kgapill{6-1}\hspace{0.25em}\allowbreak\kgapill{8-4} & {\R{流体}{りゅう|たい}・\\\R{熱}{ねつ}} \\
\bkrouteicon{2} & ②\,たし\R{合}{あ}わせ & {かけ\R{算}{ざん}→\R{面積}{めん|せき}\\→\R{積分}{せき|ぶん}} & \kgapill{1-21}\hspace{0.25em}\allowbreak\kgapill{2-3}\hspace{0.25em}\allowbreak\kgapill{4-5}\hspace{0.25em}\allowbreak\kgapill{5-10}\hspace{0.25em}\allowbreak\kgapill{6-3}\hspace{0.25em}\allowbreak\kgapill{8-11} & \R{電磁気}{でん|じ|き} \\
\bkrouteicon{3} & ③\,\R{倍々}{ばい|ばい} & {\R{累乗}{るい|じょう}→\R{指数}{し|すう}\\→$e$} & \kgapill{2-7}\hspace{0.25em}\allowbreak\kgapill{3-11}\hspace{0.25em}\allowbreak\kgapill{5-14}\hspace{0.25em}\allowbreak\kgapill{6-4}\hspace{0.25em}\allowbreak\kgapill{9-5} & \R{制御}{せい|ぎょ} \\
\bkrouteicon{4} & ④\,\R{回転}{かい|てん}と\R{波}{なみ} & {\R{角}{かく}→sin\\→\R{複素数}{ふく|そ|すう}} & \kgapill{1-20}\hspace{0.25em}\allowbreak\kgapill{4-8}\hspace{0.25em}\allowbreak\kgapill{5-15}\hspace{0.25em}\allowbreak\kgapill{6-10}\hspace{0.25em}\allowbreak\kgapill{9-1} & {\R{音}{おと}・\\\R{交流}{こう|りゅう}} \\
\bkrouteicon{5} & ⑤\,\R{表}{ひょう} & {\R{九九}{く|く}の\R{表}{ひょう}\\→\R{行列}{ぎょう|れつ}} & \kgapill{1-13}\hspace{0.25em}\allowbreak\kgapill{2-18}\hspace{0.25em}\allowbreak\kgapill{3-6}\hspace{0.25em}\allowbreak\kgapill{7-1}\hspace{0.25em}\allowbreak\kgapill{9-10} & {\R{回路網}{かい|ろ|もう}・\\\R{制御}{せい|ぎょ}} \\
\bkrouteicon{6} & ⑥\,\R{矢印}{や|じるし} & {\R{向}{む}き→ベクトル\\→\R{場}{ば}} & \kgapill{2-2}\hspace{0.25em}\allowbreak\kgapill{3-2}\hspace{0.25em}\allowbreak\kgapill{4-10}\hspace{0.25em}\allowbreak\kgapill{7-7}\hspace{0.25em}\allowbreak\kgapill{8-5} & {ロボット・\\\R{電磁気}{でん|じ|き}} \\
\bkrouteicon{7} & ⑦\,はかる & {\R{目}{め}もり→\R{誤差}{ご|さ}\\→\R{統計}{とう|けい}} & \kgapill{1-11}\hspace{0.25em}\allowbreak\kgapill{2-6}\hspace{0.25em}\allowbreak\kgapill{3-14}\hspace{0.25em}\allowbreak\kgapill{9-4} & {\R{計測}{けい|そく}・\\\R{診断}{しん|だん}} \\
\end{bktable}
\bktabcaption{7\R{本}{ほん}の\R{路線}{ろ|せん}。\R{終点}{しゅう|てん}は、\R{大学}{だい|がく}の\R{工学}{こう|がく}で\R{使}{つか}われている\R{分野}{ぶん|や}}\label{kga-routetab}
\end{bkcenter}

\clearpage

\kgaitem{\R{路線}{ろ|せん}には\R{色}{いろ}がありますが、\R{形}{かたち}の\R{記号}{き|ごう}（△・＋・×2・\R{波}{なみ}・\R{表}{ひょう}・\R{矢印}{や|じるし}・ものさし）と\R{番号}{ばん|ごう}もかならず\R{添}{そ}えてあります。\R{第}{だい}1\R{巻}{かん}では「たしあわせ」「やじるし」とひらがなで\R{書}{か}きます。}
\kgaitem{①〜④の4\R{本}{ほん}は、\R{第}{だい}6\R{巻}{かん}の\R{微分方程式}{び|ぶん|ほう|てい|しき}で\R{乗}{の}りかえます。\R{第}{だい}6\R{巻}{かん}は、シリーズ\R{全体}{ぜん|たい}の\R{乗}{の}りかえ\R{駅}{えき}です。}
\kgaitem{\R{各巻}{かく|かん}の\R{前}{まえ}のほう（\R{第}{だい}1\R{巻}{かん}は6・7ページ、ほかの\R{巻}{かん}は4・5ページ）に、\R{路線図}{ろ|せん|ず}があります。\R{第}{だい}9\R{巻}{かん}の\R{終章}{しゅう|しょう}（9-13）では、7\R{本}{ほん}の\R{路線}{ろ|せん}がどこで\R{合流}{ごう|りゅう}したかをたどります。}
\kgaitem{\R{章}{しょう}のはじめの「\R{道}{みち}しるべ」で、いまどの\R{路線}{ろ|せん}のどの\R{駅}{えき}にいるかを\R{確}{たし}かめよう。}

\bksubsection{1}{\R{路線}{ろ|せん}をたどってみよう}

\begin{bkfigure}
\begin{tikzpicture}[x=1mm, y=1mm,
  sl/.style={anchor=north, text width=21mm, align=flush center, font=\sffamily\footnotesize, text=bkInk, inner sep=0pt}]
  % ① 変化
  \node[inner sep=0pt] at (3,0) {\bkrouteicon[6mm]{1}};
  \draw[bkRoute1, line width=3pt, line cap=round] (9,0) -- (139,0);
  \foreach \c/\t [count=\i from 0] in {1-5/{ひき\R{算}{ざん}は\\かわった\R{量}{りょう}}, 2-3/{\R{変化}{へん|か}の\R{量}{りょう} Δ と\\1あたりの\R{量}{りょう}}, 3-5/{グラフの\R{傾}{かたむ}きは\\\R{変化}{へん|か}の\R{速}{はや}さ}, 5-1/{\R{平均}{へい|きん}の\R{速}{はや}さと\\\R{瞬間}{しゅん|かん}の\R{速}{はや}さ}, 6-1/{\R{微分方程式}{び|ぶん|ほう|てい|しき}って\\なに？}, 8-4/{いちばん\R{急}{きゅう}な\\\R{向}{む}き（\R{勾配}{こう|ばい}）}} {
    \fill[white] ({14+\i*23},0) circle (2.6); \draw[bkRoute1, line width=1.3pt] ({14+\i*23},0) circle (2.6);
    \node[anchor=south, inner sep=0pt] at ({14+\i*23},3.6) {\kgapill{\c}};
    \node[sl] at ({14+\i*23},-4) {\t};}
  % ⑤ 表
  \begin{scope}[shift={(0,-38)}]
  \node[inner sep=0pt] at (3,0) {\bkrouteicon[6mm]{5}};
  \draw[bkRoute5, line width=3pt, line cap=round] (9,0) -- (139,0);
  \foreach \c/\t [count=\i from 0] in {1-13/{\R{九九}{く|く}の\R{表}{ひょう}を\\たんけん}, 2-18/{\R{表計算}{ひょう|けい|さん}で\\\R{計算}{けい|さん}しよう}, 3-6/{\R{連立方程式}{れん|りつ|ほう|てい|しき}\\（\R{係数}{けい|すう}を\R{表}{ひょう}に）}, 7-1/{\R{数}{かず}の\R{表}{ひょう}で\\\R{計算}{けい|さん}する：\R{行列}{ぎょう|れつ}}, 9-10/{\R{状態方程式}{じょう|たい|ほう|てい|しき}\\（\R{行列}{ぎょう|れつ}と\R{微分方程式}{び|ぶん|ほう|てい|しき}）}} {
    \fill[white] ({16+\i*27},0) circle (2.6); \draw[bkRoute5, line width=1.3pt] ({16+\i*27},0) circle (2.6);
    \node[anchor=south, inner sep=0pt] at ({16+\i*27},3.6) {\kgapill{\c}};
    \node[sl, text width=25mm] at ({16+\i*27},-4) {\t};}
  \end{scope}
\end{tikzpicture}
\bkfigcaption{\R{変化}{へん|か}の\R{路線}{ろ|せん}（\R{上}{うえ}）と\R{表}{ひょう}の\R{路線}{ろ|せん}（\R{下}{した}）のおもな\R{駅}{えき}}
\end{bkfigure}

\R{同}{おな}じ\R{路線}{ろ|せん}の\R{駅}{えき}は、\R{前}{まえ}の\R{駅}{えき}の\R{考}{かんが}えを\R{使}{つか}って、1\R{歩}{ぽ}\R{先}{さき}へ\R{進}{すす}みます。\R{第}{だい}1\R{巻}{かん}の\R{九九}{く|く}の\R{表}{ひょう}で\R{見}{み}つけた「\R{行}{ぎょう}と\R{列}{れつ}」の\R{見方}{み|かた}は、\R{第}{だい}7\R{巻}{かん}で\R{行列}{ぎょう|れつ}という\R{大学}{だい|がく}の\R{数学}{すう|がく}になります。

\clearpage
\bksubsection{2}{\R{乗}{の}りかえ\R{駅}{えき}と\R{合流}{ごう|りゅう}}

\begin{bkfigure}
\begin{tikzpicture}[x=1mm, y=1mm]
  \foreach \r/\x/\c in {1/14/6-1, 2/44/6-3, 3/74/6-4, 4/104/6-10} {
    \node[inner sep=0pt] at (\x,58) {\bkrouteicon[6mm]{\r}};
    \draw[bkRoute\r, line width=3pt, line cap=round] (\x,53) -- (\x,0);
    \fill[white] (\x,26) circle (2.6); \draw[bkRoute\r, line width=1.3pt] (\x,26) circle (2.6);
    \node[anchor=west, inner sep=0pt] at (\x+3.6,26) {\kgapill{\c}};}
  \draw[bkInk!60, line width=0.8pt, dash pattern=on 2pt off 1.5pt, rounded corners=2.5mm] (4,16) rectangle (126,36);
  \node[anchor=north west, fill=white, inner xsep=1mm, font=\sffamily\bfseries\footnotesize, text=bkInk] at (6,37.6) {\R{乗}{の}りかえ\R{駅}{えき}：\R{第}{だい}6\R{巻}{かん}（\R{微分方程式}{び|ぶん|ほう|てい|しき}）};
\end{tikzpicture}
\bkfigcaption{①〜④の4\R{本}{ほん}の\R{路線}{ろ|せん}は、\R{第}{だい}6\R{巻}{かん}の\R{微分方程式}{び|ぶん|ほう|てい|しき}で\R{乗}{の}りかえる}
\end{bkfigure}

\R{第}{だい}6\R{巻}{かん}は「\R{次}{つぎ}の\R{瞬間}{しゅん|かん}を\R{決}{き}めるルール」（\R{微分方程式}{び|ぶん|ほう|てい|しき}）の\R{巻}{かん}です。\R{変化}{へん|か}・たし\R{合}{あ}わせ・\R{倍々}{ばい|ばい}・\R{回転}{かい|てん}と\R{波}{なみ}の4つの\R{考}{かんが}えが、ここで\R{出}{で}あいます。\R{第}{だい}9\R{巻}{かん}の\R{終章}{しゅう|しょう}（9-13）では、\R{路線}{ろ|せん}どうしが\R{合流}{ごう|りゅう}する5つの\R{所}{ところ}をたどります。いまはわからなくてもかまいません。\R{名前}{な|まえ}だけ\R{覚}{おぼ}えておいて、\R{第}{だい}9\R{巻}{かん}で\R{確}{たし}かめよう。

\begin{bkcenter}
\begin{bktable}{colspec={Q[c,m,wd=8mm] Q[l,m,wd=72mm] Q[l,m,wd=36mm]}, rowsep=1.6pt, cells={font=\small}, row{1}={font=\sffamily\bfseries\small}}
 & 9-13 で\R{合流}{ごう|りゅう}する\R{所}{ところ} & \R{合流}{ごう|りゅう}する\R{路線}{ろ|せん} \\
（ア） & $s=\sigma+j\omega$：\R{倍々}{ばい|ばい}と\R{回転}{かい|てん}が1つの\R{数}{かず}に & \bkrouteicon{3}\,\bkrouteicon{4} \\
（イ） & $e^{At}$：\R{倍々}{ばい|ばい}と\R{表}{ひょう}（\R{行列}{ぎょう|れつ}） & \bkrouteicon{3}\,\bkrouteicon{5} \\
（ウ） & \R{熱}{ねつ}の\R{行列}{ぎょう|れつ}の\R{固有}{こ|ゆう}ベクトルは sin の\R{形}{かたち} & \bkrouteicon{1}\,\bkrouteicon{5}\,\bkrouteicon{4} \\
（エ） & ガウスの\R{形}{かたち} & \bkrouteicon{2}\,\bkrouteicon{7}\,\bkrouteicon{4} \\
（オ） & \R{波動方程式}{は|どう|ほう|てい|しき}と\R{音}{おと}・\R{光}{ひかり} & \bkrouteicon{1}\,\bkrouteicon{6}\,\bkrouteicon{4} \\
\end{bktable}
\bktabcaption{\R{終章}{しゅう|しょう} 9-13 でたどる5つの\R{合流}{ごう|りゅう}}
\end{bkcenter}

\vspace{2mm}
\noindent{\small \R{次}{つぎ}のページの\R{路線図}{ろ|せん|ず}の\R{駅}{えき}の\R{番号}{ばん|ごう}は\R{章}{しょう}の\R{番号}{ばん|ごう}です（「5-1」は\R{第}{だい}5\R{巻}{かん}の1\R{章}{しょう}）。\R{番号}{ばん|ごう}の\R{字}{じ}が\R{小}{ちい}さくて\R{読}{よ}みにくいときは、\R{表}{ひょう}1（\bkpageno{\pageref{kga-routetab}}）で\R{確}{たし}かめよう。}\par


\clearpage
\noindent{\sffamily\bfseries\fontsize{20}{26}\selectfont\color{\bkTT}\R{全}{ぜん}9\R{巻}{かん}の\R{路線図}{ろ|せん|ず}}\hfill{\sffamily\fontsize{8.5}{11}\selectfont\color{bkGray}\R{下}{した}へ\R{行}{い}くほど\R{先}{さき}の\R{巻}{かん}。○が\R{駅}{えき}（\R{章}{しょう}）です。}\par\vspace{1mm}
{\color{bkTheme}\hrule height 1.4pt}\vspace{2mm}
\noindent\bkSeriesMap\par


%% =====================================================================
%% 4　1つの章の歩き方（4ページ）
%% =====================================================================
\kgasec{4}{1つの\R{章}{しょう}の\R{歩}{ある}き\R{方}{かた}}
  {どの\R{章}{しょう}も、いつも\R{同}{おな}じ\R{部品}{ぶ|ひん}が、\R{同}{おな}じ\R{順番}{じゅん|ばん}でならんでいます。\R{部品}{ぶ|ひん}の\R{名前}{な|まえ}と\R{役目}{やく|め}を\R{知}{し}っておくと、どの\R{巻}{かん}でも\R{迷}{まよ}わずに\R{読}{よ}めます。}
  {1つの章の歩き方}

\R{下}{した}の\R{図}{ず}は、1つの\R{章}{しょう}のページの\R{見本}{み|ほん}です。\R{番号}{ばん|ごう}は\R{次}{つぎ}のページからの\R{表}{ひょう}と\R{同}{おな}じです。\R{章}{しょう}は、いつも\R{見開}{み|ひら}きの\R{左}{ひだり}のページから\R{始}{はじ}まります。

\begin{bkfigure}
\begin{tikzpicture}[x=0.93mm, y=0.7mm]
  % ---- 見開き1（章のはじめ） ----
  \foreach \px in {0,47} {\fill[bkLine!50] (\px+0.8,-0.8) rectangle ++(46,70); \draw[bkGray, line width=0.6pt, fill=white] (\px,0) rectangle ++(46,70);}
  % 左ページ
  \fill[bkVol3, rounded corners=1mm] (3,57) rectangle (43,67);
  \fill[white] (8,62) circle (3.3); \node[font=\sffamily\bfseries\footnotesize, text=bkOn3] at (8,62) {3-5};
  \fill[white] (13,63.5) rectangle (38,64.6); \fill[white] (13,60.4) rectangle (32,61.5);
  \draw[bkGray, line width=0.4pt] (3,54.4) -- (43,54.4); \fill[bkVol3] (16,55) rectangle ++(5,1.8); \fill[bkVol2] (22,55) rectangle ++(5,1.8);
  \draw[bkVol3!70!black, line width=0.7pt, rounded corners=1mm] (3,43) rectangle (43,52.5);
  \foreach \y in {49.5,47,44.5} \fill[bkStar] (5,\y) rectangle ++(3,1); \foreach \y in {49.5,47,44.5} \fill[bkLine] (10,\y) rectangle ++(28,1);
  \draw[bkYellow, line width=0.7pt, fill=bkYellowBg, rounded corners=1mm] (3,31) rectangle (43,41);
  \foreach \y in {36,33.5} \fill[bkLine] (6,\y) rectangle ++(33,1);
  \draw[bkGreen, line width=0.7pt, fill=bkGreenBg, rounded corners=1mm] (3,15) rectangle (43,29);
  \foreach \y in {24,21.5,19,16.5} \fill[bkLine] (6,\y) rectangle ++(33,1);
  \draw[bkGray!60, line width=0.6pt, fill=bkPaper, rounded corners=1mm] (3,3) rectangle (43,12.5);
  \foreach \y in {8,5.5} \fill[bkLine] (6,\y) rectangle ++(33,1);
  % 右ページ
  \fill[bkVol3, rounded corners=0.5mm] (50,63) circle (2); \fill[bkLine] (54,62.4) rectangle ++(25,1.4);
  \foreach \y in {58.5,56,53.5,51} \fill[bkLine] (50,\y) rectangle ++(40,1);
  \draw[bkGray, line width=0.6pt, dashed, fill=bkPaper, rounded corners=1mm] (50,38) rectangle (90,48);
  \foreach \y in {43,40.5} \fill[bkLine] (53,\y) rectangle ++(33,1);
  \fill[bkNavy] (50,26) rectangle ++(1,9); \foreach \y in {32,29.5,27} \fill[bkLine] (53,\y) rectangle ++(33,1);
  \draw[bkNavy, line width=0.9pt, rounded corners=1mm] (50,14) rectangle (90,23); \fill[bkLine] (54,18) rectangle ++(32,1.2);
  \draw[bkVol3!70!black, line width=0.7pt, rounded corners=1mm] (50,2.5) rectangle (68,11.5); \fill[bkVol3] (50.35,9) rectangle (67.65,11.15);
  \node[font=\footnotesize, text=bkInk] at (55,6) {$\Delta$};
  \foreach \y in {8,5.5} \fill[bkLine] (71,\y) rectangle ++(19,1);
  % 番号
  \kganum{(-3.6,64.5)}{1}\kganum{(-3.6,54.4)}{2}\kganum{(-3.6,47.8)}{3}\kganum{(-3.6,36)}{4}\kganum{(-3.6,22)}{5}\kganum{(-3.6,7.8)}{6}
  \kganum{(95.6,57)}{7}\kganum{(95.6,43)}{8}\kganum{(95.6,30.5)}{9}\kganum{(95.6,18.5)}{10}\kganum{(95.6,7)}{11}
  \node[anchor=south, font=\sffamily\bfseries\footnotesize, text=bkGray] at (46.5,70.6) {\R{見開}{み|ひら}き\,A\enspace（\R{章}{しょう}のはじめ）};
  % ---- 見開き2（章の終わり） ----
  \begin{scope}[shift={(0,-84)}]
  \foreach \px in {0,47} {\fill[bkLine!50] (\px+0.8,-0.8) rectangle ++(46,70); \draw[bkGray, line width=0.6pt, fill=white] (\px,0) rectangle ++(46,70);}
  \draw[bkBlue, line width=0.7pt, fill=bkBlueBg, rounded corners=1mm] (3,48) rectangle (43,67);
  \foreach \y in {61,58.5,56,53.5,51} \fill[bkLine] (6,\y) rectangle ++(33,1);
  \draw[bkOrange, line width=1pt, rounded corners=1mm] (3,35) rectangle (43,45);
  \foreach \y in {40,37.5} \fill[bkLine] (6,\y) rectangle ++(33,1);
  \draw[bkRed, line width=0.7pt, fill=bkRedBg, rounded corners=1mm] (3,20) rectangle (43,32);
  \foreach \y in {27,24.5,22} \fill[bkLine] (6,\y) rectangle ++(33,1);
  \draw[bkPurple, line width=0.7pt, fill=bkPurpleBg, rounded corners=1mm] (3,3) rectangle (43,17);
  \foreach \y in {12,9.5,7,4.5} \fill[bkLine] (6,\y) rectangle ++(33,1);
  \fill[bkVol3, rounded corners=0.5mm] (50,57) rectangle (90,66);
  \foreach \y in {52,49.5,47,44.5,42,39.5,37,34.5} \fill[bkLine] (50,\y) rectangle ++(40,1);
  \draw[bkVol3!80!black, line width=0.7pt, fill=bkVol3!8, rounded corners=1mm] (50,3) rectangle (90,23);
  \foreach \y in {18,14.5,11,7.5} {\draw[bkVol3!80!black, line width=0.5pt] (52.5,\y) rectangle ++(2,2); \fill[bkLine] (56,\y+0.5) rectangle ++(31,1);}
  \kganum{(-3.6,57.5)}{12}\kganum{(-3.6,40)}{13}\kganum{(-3.6,26)}{14}\kganum{(-3.6,10)}{15}
  \kganum{(95.6,48)}{16}\kganum{(95.6,13)}{17}
  \node[anchor=south, font=\sffamily\bfseries\footnotesize, text=bkGray] at (46.5,70.6) {\R{見開}{み|ひら}き\,B\enspace（\R{章}{しょう}の\R{終}{お}わり）};
  % 巻末のこたえ
  \fill[bkLine!50] (104.8,9.2) rectangle ++(32,50); \draw[bkGray, line width=0.6pt, fill=white] (104,10) rectangle ++(32,50);
  \node[anchor=north west, font=\sffamily\bfseries\footnotesize, text=bkInk, inner sep=0pt] at (106.5,57.5) {こたえ};
  \draw[bkTheme, line width=0.8pt] (106.5,52.8) -- (133.5,52.8);
  \foreach \y in {48,45,42,39,36,33,30,27,24,21,18,15} \fill[bkLine] (106.5,\y) rectangle ++(12,1);
  \foreach \y in {48,45,42,39,36,33,30,27,24,21,18,15} \fill[bkLine] (121,\y) rectangle ++(12,1);
  \kganum{(141,35)}{18}
  \node[anchor=south, font=\sffamily\bfseries\footnotesize, text=bkGray] at (120,60.6) {\R{本}{ほん}の\R{後}{うし}ろ};
  \end{scope}
  % 右上のあき：読み方
  \node[anchor=north west, text width=40mm, align=left, font=\sffamily\footnotesize, text=bkInk, inner sep=0pt] at (102,66)
    {\R{見本}{み|ほん}は\R{第}{だい}3\R{巻}{かん}の\R{色}{いろ}でかきました。\R{本}{ほん}では、1つの\R{部品}{ぶ|ひん}が\R{何回}{なん|かい}も\R{出}{で}てくることもあります。\par\vspace{1.5mm}
     \R{部品}{ぶ|ひん}の\R{色}{いろ}は、どの\R{巻}{かん}でも\R{同}{おな}じです。\R{色}{いろ}のほかに、\R{部品}{ぶ|ひん}ごとのマークと\R{名前}{な|まえ}でも\R{見分}{み|わ}けられます。};
\end{tikzpicture}
\bkfigcaption{1つの\R{章}{しょう}のページの\R{見本}{み|ほん}（\R{番号}{ばん|ごう}は\R{表}{ひょう}の①〜⑱）}
\end{bkfigure}

\clearpage
\begin{bkcenter}
\begin{bktable}{colspec={Q[c,m,wd=7mm] Q[c,m,wd=8mm] Q[l,m,wd=30mm] Q[l,m,wd=77mm]}, rowsep=0.6pt, cells={font=\small}, row{1}={font=\sffamily\bfseries\small}}
 & & \R{部品}{ぶ|ひん} & \R{役目}{やく|め}と\R{使}{つか}い\R{方}{かた} \\
① & \tikz\fill[bkTheme, rounded corners=0.6mm] (0,0) rectangle (5mm,3.2mm); & \R{章}{しょう}の\R{帯}{おび} & \R{章}{しょう}の\R{番号}{ばん|ごう}・\R{題名}{だい|めい}・★。\R{題名}{だい|めい}は、みんなに\R{問}{と}いかける\R{形}{かたち}が\R{多}{おお}い。 \\
② & \tikz{\fill[bkTheme] (0,0) circle (1.6mm); \draw[white, line width=0.6pt] (0,-1.1mm)--(0,1.1mm);} & \R{道}{みち}しるべ & 「この\R{章}{しょう}の\R{前}{まえ}に」\R{読}{よ}むとよい\R{章}{しょう}と、「この\R{章}{しょう}の\R{先}{さき}に」この\R{考}{かんが}えが\R{育}{そだ}っていく\R{章}{しょう}。 \\
③ & \bkicon[1.3em]{goal}{bkTheme} & この\R{章}{しょう}でできるようになること & この\R{章}{しょう}のゴール（ふつう3つ）。1つずつに★がついています。 \\
④ & \bkicon[1.3em]{hatena}{bkYellow!80!black} & はてな？ & \R{身近}{み|ぢか}なふしぎ（\R{章}{しょう}の\R{問}{と}い）。\R{答}{こた}えは、かならずこの\R{章}{しょう}の\R{中}{なか}にあります。 \\
⑤ & \bkicon[1.3em]{yattemiyou}{bkGreen} & やってみよう & \R{実験}{じっ|けん}・\R{観察}{かん|さつ}・\R{電卓}{でん|たく}・\R{表計算}{ひょう|けい|さん}。\R{手順}{て|じゅん}に\R{番号}{ばん|ごう}があります。 \\
⑥ & \bkicon[1.3em]{omoidasu}{bkGray} & \R{思}{おも}い\R{出}{だ}そう & \R{前}{まえ}の\R{巻}{かん}・\R{前}{まえ}の\R{章}{しょう}の\R{大事}{だい|じ}なところを2〜4\R{行}{ぎょう}で。\R{章}{しょう}の\R{番号}{ばん|ごう}つきなので、もどって\R{確}{たし}かめられます。 \\
⑦ & \bkicon[1.3em]{doguu}{bkNavy} & \R{数学}{すう|がく}の\R{道具}{どう|ぐ} & \R{章}{しょう}の\R{中心}{ちゅう|しん}。\R{小見出}{こ|み|だ}し（①②③…）ごとに1つの\R{考}{かんが}え。 \\
⑧ & \bkicon[1.3em]{motto}{bkGray} & もっと\R{知}{し}りたい & もう\R{一歩}{いっ|ぽ}\R{先}{さき}の\R{話}{はなし}。★つき。とばして\R{先}{さき}へ\R{進}{すす}んでもかまいません。 \\
⑨ & \tikz\fill[bkNavy] (0,0) rectangle (1mm,4mm); & せいかくに\R{言}{い}うと & たとえ\R{話}{ばなし}のすぐあと。どこまで\R{正}{ただ}しいかと、\R{正確}{せい|かく}な\R{決}{き}まり。 \\
\end{bktable}
\bktabcaption{\R{章}{しょう}の\R{部品}{ぶ|ひん}（1）\R{章}{しょう}のはじめから\R{数学}{すう|がく}の\R{道具}{どう|ぐ}まで}
\end{bkcenter}

\bksubsection{1}{\R{道}{みち}しるべの\R{読}{よ}み\R{方}{かた}}

\R{章}{しょう}の\R{帯}{おび}のすぐ\R{下}{した}に、こんな\R{行}{ぎょう}があります（\R{見本}{み|ほん}）。
\par\vspace{1mm}
\noindent\begin{minipage}{\linewidth}\bkmichishirube{1-14,2-2}{3-5,5-9}\end{minipage}
\R{札}{ふだ}の\R{色}{いろ}は、その\R{章}{しょう}がある\R{巻}{かん}の\R{色}{いろ}です。わからないことばや\R{式}{しき}が\R{出}{で}てきたら、「この\R{章}{しょう}の\R{前}{まえ}に」の\R{章}{しょう}にもどって\R{確}{たし}かめよう。\R{前}{まえ}の\R{巻}{かん}にあることは、\R{思}{おも}い\R{出}{だ}そう（⑥）で\R{短}{みじか}くふり\R{返}{かえ}るだけで、\R{同}{おな}じ\R{説明}{せつ|めい}はくり\R{返}{かえ}しません。

\clearpage

\begin{bkcenter}
\begin{bktable}{colspec={Q[c,m,wd=7mm] Q[c,m,wd=8mm] Q[l,m,wd=30mm] Q[l,m,wd=77mm]}, rowsep=0.6pt, cells={font=\small}, row{1}={font=\sffamily\bfseries\small}}
 & & \R{部品}{ぶ|ひん} & \R{役目}{やく|め}と\R{使}{つか}い\R{方}{かた} \\
⑩ & \bkicon[1.3em]{taisetsu}{bkNavy} & たいせつ & いちばん\R{大事}{だい|じ}なこと。1\R{行}{ぎょう}に1つ。 \\
⑪ & \bkicon[1.3em]{card}{bkTheme} & \R{記号}{き|ごう}カード & \R{新}{あたら}しい\R{記号}{き|ごう}の\R{読}{よ}み\R{方}{かた}・\R{意味}{い|み}・\R{英語}{えい|ご}。1つの\R{章}{しょう}に2\R{枚}{まい}まで。\R{巻末}{かん|まつ}にも\R{一覧}{いち|らん}があります。 \\
⑫ & \bkicon[1.3em]{reidai}{bkBlue} & \R{例題}{れい|だい} & \R{解}{と}き\R{方}{かた}のお\R{手本}{て|ほん}。「\R{考}{かんが}え\R{方}{かた}」「しき」「こたえ」の\R{札}{ふだ}（\R{第}{だい}4\R{巻}{かん}からは「\R{式}{しき}」「\R{答}{こた}え」）で、\R{途中}{と|ちゅう}の\R{式}{しき}も1\R{行}{ぎょう}ずつ\R{書}{か}いてあります。 \\
⑬ & \bkicon[1.3em]{jibun}{bkOrange} & じぶんでとこう & \R{例題}{れい|だい}と\R{同}{おな}じ\R{形}{かたち}の\R{問題}{もん|だい}。お\R{手本}{て|ほん}をまねして、ノートに\R{解}{と}こう。 \\
⑭ & \bkicon[1.3em]{tsumazuki}{bkRed} & まちがえやすいところ & ×と○を1\R{行}{ぎょう}ずつならべ、「なぜまちがいか」を\R{書}{か}いています。 \\
⑮ & \bkicon[1.3em]{butsuri}{bkPurple} & \R{物理}{ぶつ|り}・\R{工学}{こう|がく}ではこう\R{使}{つか}う & \R{本物}{ほん|もの}の\R{機械}{き|かい}やしくみで、この\R{章}{しょう}の\R{数学}{すう|がく}がどう\R{働}{はたら}いているか。 \\
⑯ & \bkicon[1.3em]{check}{bkTheme} & \R{章末問題}{しょう|まつ|もん|だい} & 「\R{基本}{き|ほん}」と「チャレンジ」。\R{問題}{もん|だい}の\R{形}{かたち}は、どれも\R{例題}{れい|だい}かじぶんでとこうで1\R{回}{かい}\R{見}{み}せてあります。 \\
⑰ & \bkicon[1.3em]{check}{bkTheme} & ふりかえり & ③と\R{同}{おな}じゴールを、チェックの\R{形}{かたち}で。「→〇ページ」で、その\R{説明}{せつ|めい}にもどれます。 \\
⑱ & \bkicon[1.3em]{reidai}{bkGray} & こたえ（\R{本}{ほん}の\R{後}{うし}ろ） & じぶんでとこうと\R{章末問題}{しょう|まつ|もん|だい}のこたえ。\R{使}{つか}い\R{方}{かた}は8\R{節}{せつ}。 \\
\end{bktable}
\bktabcaption{\R{章}{しょう}の\R{部品}{ぶ|ひん}（2）まとめから\R{章}{しょう}の\R{終}{お}わりまで}
\end{bkcenter}

\bksubsection{2}{\R{第}{だい}1\R{巻}{かん}だけのちがい}

\begin{bkcenter}
\begin{bktable}{colspec={Q[l,m,wd=46mm] Q[l,m,wd=46mm] Q[l,m,wd=30mm]}, rowsep=1.4pt, cells={font=\small}, row{1}={font=\sffamily\bfseries\small}}
\R{第}{だい}2\R{巻}{かん}から & \R{第}{だい}1\R{巻}{かん} & \R{部品}{ぶ|ひん} \\
\R{例題}{れい|だい} & れいだい & ⑫ \\
\R{章末問題}{しょう|まつ|もん|だい} & もんだい & ⑯ \\
\R{物理}{ぶつ|り}・\R{工学}{こう|がく}ではこう\R{使}{つか}う & つかわれている ところ & ⑮ \\
\end{bktable}
\bktabcaption{\R{第}{だい}1\R{巻}{かん}のことばと\R{使}{つか}い\R{方}{かた}}
\end{bkcenter}



\bksubsection{3}{\R{部品}{ぶ|ひん}の\R{順番}{じゅん|ばん}}

\R{部品}{ぶ|ひん}は、だいたい\R{次}{つぎ}の\R{順番}{じゅん|ばん}でならんでいます。\R{章}{しょう}の\R{長}{なが}さは\R{章}{しょう}ごとにちがい、14ページの\R{章}{しょう}も、30ページをこえる\R{章}{しょう}もあります。

\begin{bkfigure}
\begin{tikzpicture}[x=1mm, y=1mm,
  st/.style={draw=#1, line width=0.9pt, fill=white, rounded corners=1.8mm, minimum width=31mm, minimum height=11mm, align=flush center, font=\sffamily\footnotesize, text=bkInk, inner sep=0.8mm}]
  \foreach \n/\t/\c [count=\i from 0] in {1/{\R{章}{しょう}の\R{帯}{おび}・\R{道}{みち}しるべ}/bkTheme, 3/{できるようになること}/bkTheme, 4/{はてな？}/bkYellow!85!black, 5/{やってみよう}/bkGreen,
      6/{\R{思}{おも}い\R{出}{だ}そう}/bkGray, 7/{\R{数学}{すう|がく}の\R{道具}{どう|ぐ}\\（⑧⑨⑩⑪をふくむ）}/bkNavy, 12/{\R{例題}{れい|だい}と\\じぶんでとこう}/bkBlue, 14/{まちがえやすい\\ところ}/bkRed,
      15/{\R{物理}{ぶつ|り}・\R{工学}{こう|がく}では\\こう\R{使}{つか}う}/bkPurple, 16/{\R{章末問題}{しょう|まつ|もん|だい}}/bkTheme, 17/{ふりかえり}/bkTheme, 18/{こたえ（\R{巻末}{かん|まつ}）}/bkGray} {
    \pgfmathtruncatemacro\kgar{\i/4}\pgfmathtruncatemacro\kgac{mod(\i,4)}
    \node[st=\c] (f\i) at ({\kgac*35.6+15.5},{-\kgar*19}) {\t};
    \node[circle, fill=bkInk, text=white, inner sep=0pt, minimum size=4.4mm, font=\sffamily\bfseries\footnotesize] at (f\i.north west) {\n};}
  \foreach \a/\b in {0/1,1/2,2/3,4/5,5/6,6/7,8/9,9/10,10/11} \draw[-{Stealth[length=2mm]}, line width=1pt, bkGray] ([xshift=0.6mm]f\a.east) -- ([xshift=-0.6mm]f\b.west);
  \draw[-{Stealth[length=2mm]}, line width=1pt, bkGray, rounded corners=2mm] ([yshift=-0.6mm]f3.south) -- ++(0,-2.6) -| ([yshift=0.6mm]f4.north);
  \draw[-{Stealth[length=2mm]}, line width=1pt, bkGray, rounded corners=2mm] ([yshift=-0.6mm]f7.south) -- ++(0,-2.6) -| ([yshift=0.6mm]f8.north);
\end{tikzpicture}
\bkfigcaption{1つの\R{章}{しょう}の\R{流}{なが}れ（\R{番号}{ばん|ごう}は\R{表}{ひょう}の①〜⑱）}
\end{bkfigure}

\begin{yattemiyou}[\R{本物}{ほん|もの}の\R{章}{しょう}で、\R{部品}{ぶ|ひん}をさがそう]
\begin{enumerate}[label={\bkcircled[bkGreen]{\arabic*}}, leftmargin=2.4\zw, itemsep=1pt, topsep=2pt]
\item どれか1\R{冊}{さつ}を\R{開}{ひら}いて、もくじから\R{好}{す}きな\R{章}{しょう}を1つ\R{選}{えら}ぼう。
\item ページをめくりながら①〜⑰の\R{部品}{ぶ|ひん}をさがして、ノートに\R{番号}{ばん|ごう}とページを\R{書}{か}こう。
\item 「はてな？」の\R{問}{と}いの\R{答}{こた}えが、\R{章}{しょう}のどこに\R{書}{か}いてあるかを\R{見}{み}つけよう。
\end{enumerate}
\end{yattemiyou}

\bkfukidashi{mina}{\R{章}{しょう}の\R{最後}{さい|ご}の「ふりかえり」で、できなかった\R{項目}{こう|もく}があったら、「→〇ページ」のページにもどってみよう。}


%% =====================================================================
%% 5　★と、3つの段（4ページ）
%% =====================================================================
\kgasec{5}{★と、3つの\R{段}{だん}}
  {★は「どこまで\R{登}{のぼ}るか」のしるし。\R{大学}{だい|がく}の\R{内容}{ない|よう}は、「\R{目}{め}でわかる→\R{手}{て}でわかる→\R{式}{しき}でわかる」の3つの\R{段}{だん}で\R{登}{のぼ}ります。}
  {★と、3つの段}

\bksubsection{1}{★は\R{山}{やま}の\R{高}{たか}さ}

\noindent\begin{minipage}[c]{0.40\linewidth}\centering
\begin{tikzpicture}[x=1mm, y=1mm]
  \foreach \k/\c in {1/14,2/28,3/44,4/62,5/82} {
    \pgfmathsetmacro\kgayb{(\k-1)*9}\pgfmathsetmacro\kgayt{\k*9}
    \fill[bkTheme!\c] ({36-19*(1-\kgayb/45)},\kgayb) -- ({36+19*(1-\kgayb/45)},\kgayb)
      -- ({36+19*(1-\kgayt/45)},\kgayt) -- ({36-19*(1-\kgayt/45)},\kgayt) -- cycle;
    \node[anchor=east, inner sep=0pt] at (15,{(\k-0.5)*9}) {\bkstarmeter{\k}{\k}};}
  \draw[bkInk!70, line width=0.8pt] (36,45) -- (36,52);
  \fill[bkRed] (36,52) -- (42,50) -- (36,48) -- cycle;
\end{tikzpicture}
\end{minipage}\hfill
\begin{minipage}[c]{0.57\linewidth}
\begin{bktable}{colspec={Q[l,m,wd=24mm] Q[l,m,wd=44mm]}, rowsep=1pt, cells={font=\small}, row{1}={font=\sffamily\bfseries\small}}
★の\R{数}{かず} & \R{山登}{やま|のぼ}りにたとえると \\
\bkstarmeter{1}{1} & はじめの\R{一歩}{いっ|ぽ} \\
\bkstarmeter{2}{2} & ふもとを\R{歩}{ある}く \\
\bkstarmeter{3}{3} & \R{山}{やま}を\R{登}{のぼ}る \\
\bkstarmeter{4}{4} & \R{尾根}{お|ね}を\R{歩}{ある}く \\
\bkstarmeter{5}{5} & \R{頂上}{ちょう|じょう}をめざす \\
\end{bktable}
\end{minipage}
\par\vspace{2mm}

\R{章}{しょう}・「できるようになること」の1つ1つ・もっと\R{知}{し}りたいには、★がついています。★は、このシリーズの\R{中}{なか}での\R{難}{むずか}しさの\R{目安}{め|やす}で、どんな\R{道具}{どう|ぐ}（\R{計算}{けい|さん}のしかたや\R{考}{かんが}え\R{方}{かた}）を\R{使}{つか}うかで\R{決}{き}めています。\R{学校}{がっ|こう}の\R{学年}{がく|ねん}と、ぴったり\R{同}{おな}じではありません。

★が「★〜★★」のように\R{範囲}{はん|い}で\R{書}{か}いてある\R{章}{しょう}もあります。もくじでは、★（ぬりつぶし）がその\R{章}{しょう}でかならず\R{登}{のぼ}る\R{高}{たか}さ、☆（\R{白}{しろ}ぬき）が「この\R{章}{しょう}の\R{中}{なか}で、ここまで\R{登}{のぼ}る」というしるしです。たとえば \bkstarmeter{2}{3} は、「★★まではかならず、★★★まで\R{登}{のぼ}る」と\R{読}{よ}みます。\R{自分}{じ|ぶん}の\R{読}{よ}める★まで\R{読}{よ}んで、\R{先}{さき}へ\R{進}{すす}んでもだいじょうぶ。



\R{第}{だい}4\R{巻}{かん}のように、1\R{冊}{さつ}の\R{中}{なか}に★から★★★★★までの\R{章}{しょう}がある\R{巻}{かん}もあります。★の\R{高}{たか}い\R{章}{しょう}で\R{止}{と}まってしまったら、その\R{章}{しょう}の☆の\R{手前}{て|まえ}まで\R{読}{よ}んで\R{次}{つぎ}の\R{章}{しょう}へ\R{進}{すす}み、あとでもどってきてもかまいません。

\bksubsection{2}{3つの\R{段}{だん}：\R{目}{め}でわかる→\R{手}{て}でわかる→\R{式}{しき}でわかる}

★★★★★の\R{章}{しょう}（\R{大学}{だい|がく}で\R{学}{まな}ぶ\R{内容}{ない|よう}）は\R{第}{だい}4\R{巻}{かん}から\R{出}{で}てきて、\R{第}{だい}6\R{巻}{かん}から\R{先}{さき}で\R{多}{おお}くなります。\R{第}{だい}6\R{巻}{かん}から\R{先}{さき}のそうした\R{章}{しょう}の\R{多}{おお}くは、3つの\R{段}{だん}に\R{分}{わ}けて\R{書}{か}いてあります。\R{段}{だん}の★は、その\R{段}{だん}のゴールの\R{高}{たか}さです。はじめの2つの\R{段}{だん}（\R{目}{め}・\R{手}{て}）は、★★★までに\R{習}{なら}った\R{計算}{けい|さん}のしかただけで\R{進}{すす}めるように\R{書}{か}いてあります。

\begin{bkfigure}
\begin{tikzpicture}[x=1mm, y=1mm]
  \foreach \ic/\nm/\tx/\lv [count=\i from 0] in {me/{\R{目}{め}でわかる}/{\R{図}{ず}と\R{実験}{じっ|けん}で、\\\R{何}{なに}が\R{起}{お}きているかを\R{見}{み}る}/3,
                                         te/{\R{手}{て}でわかる}/{\R{表計算}{ひょう|けい|さん}やプログラムで、\\\R{数}{かず}を\R{出}{だ}して\R{確}{たし}かめる}/4,
                                         shiki/{\R{式}{しき}でわかる}/{\R{大学}{だい|がく}の\R{式}{しき}を\R{作}{つく}り、\\\R{式}{しき}で\R{確}{たし}かめる}/5} {
    \node[draw=bkTheme, line width=1pt, fill=white, rounded corners=2.5mm, minimum width=40mm, minimum height=41mm, anchor=south west] (d\i) at (\i*48,\i*7) {};
    \node[anchor=north] (di\i) at ([yshift=-2mm]d\i.north) {\bkicon[1.8em]{\ic}{bkTheme}};
    \node[anchor=north, font=\sffamily\bfseries\small, text=\bkTT, inner sep=0pt] (dn\i) at ([yshift=-1.5mm]di\i.south) {\nm};
    \node[anchor=north, text width=37mm, align=flush center, font=\sffamily\footnotesize, text=bkInk, inner sep=0pt] at ([yshift=-2mm]dn\i.south) {\tx};
    \node[anchor=south, inner sep=0pt] at ([yshift=1.5mm]d\i.south) {\bklevel{\lv}};}
  \draw[-{Stealth[length=2.4mm]}, line width=1.1pt, bkGray] ([xshift=0.8mm]d0.east) -- ([xshift=-0.8mm]d1.west |- d0.east);
  \draw[-{Stealth[length=2.4mm]}, line width=1.1pt, bkGray] ([xshift=0.8mm]d1.east) -- ([xshift=-0.8mm]d2.west |- d1.east);
  % ここまでで十分の旗
  \draw[bkInk!70, line width=0.8pt] (88,48) -- (88,61);
  \fill[bkTheme] (88,61) -- (95,58.5) -- (88,56) -- cycle;
  \node[anchor=east, font=\sffamily\bfseries\footnotesize, text=\bkTT, inner sep=1pt] at (87,58) {ここまでで\R{十分}{じゅう|ぶん}};
\end{tikzpicture}
\bkfigcaption{3つの\R{段}{だん}。「\R{手}{て}でわかる」の\R{終}{お}わりに「ここまでで\R{十分}{じゅう|ぶん}」の\R{旗}{はた}が\R{立}{た}つ。★は、\R{段}{だん}のゴールのよくある\R{高}{たか}さ（\R{章}{しょう}によってちがう）}
\end{bkfigure}

たとえば\R{第}{だい}8\R{巻}{かん}の 8-9「わき\R{出}{だ}しと\R{吸}{す}いこみ」では、\R{目}{め}でわかる「\R{小}{ちい}さな\R{箱}{はこ}を\R{置}{お}いてみる」（★★★）→\nobreak\R{手}{て}でわかる「マスの4つの\R{辺}{へん}で\R{数}{かぞ}える」（★★★★）→\nobreak\R{式}{しき}でわかる「マスを\R{小}{ちい}さくしていく」（★★★★★）の\R{順}{じゅん}に\R{進}{すす}みます。\R{段}{だん}の\R{見出}{み|だ}しには、\R{目}{め}・\R{表}{ひょう}・$\partial$ のマークがついています。

「\R{手}{て}でわかる」の\R{終}{お}わりには、\R{次}{つぎ}の\R{旗}{はた}が\R{立}{た}っています。

\bkkokomade

この\R{旗}{はた}まで\R{読}{よ}めたら、その\R{章}{しょう}のゴールです。「\R{式}{しき}でわかる」は、\R{式}{しき}で\R{確}{たし}かめたい\R{人}{ひと}や、\R{大学}{だい|がく}の\R{専門書}{せん|もん|しょ}を\R{読}{よ}む\R{準備}{じゅん|び}をしたい\R{人}{ひと}のための\R{段}{だん}です。\R{旗}{はた}のところで\R{次}{つぎ}の\R{章}{しょう}へ\R{進}{すす}んでもかまいません。

3つの\R{段}{だん}は、6-10〜6-14、7-7〜7-10、8-1〜8-13、9-1〜9-11 の\R{章}{しょう}で\R{使}{つか}っています（8-5 は2つの\R{段}{だん}）。

\bksubsection{3}{\R{第}{だい}8・9\R{巻}{かん}：「\R{名前}{な|まえ}」と「\R{実際}{じっ|さい}の\R{操作}{そう|さ}」}

ここは、\R{第}{だい}8・9\R{巻}{かん}を\R{読}{よ}むときに、もう\R{一度}{いち|ど}\R{開}{ひら}くページです。いまは\R{名前}{な|まえ}だけ\R{知}{し}っておけば、だいじょうぶ。

\R{第}{だい}8\R{巻}{かん}と\R{第}{だい}9\R{巻}{かん}には、$\textstyle\iint_D$、$\textstyle\oint_C$、$\mathcal{L}[f(t)]$、$G(s)$ のような\R{記号}{き|ごう}がたくさん\R{出}{で}てきます。これらの\R{記号}{き|ごう}は、「こういう\R{計算}{けい|さん}をする」という\R{操作}{そう|さ}の\textbf{\R{名前}{な|まえ}}です。\R{記号}{き|ごう}を\R{見}{み}ただけでは、まだ\R{計算}{けい|さん}はできません。\R{実際}{じっ|さい}の\R{値}{あたい}は、\R{知}{し}っている\R{計算}{けい|さん}に\textbf{\R{直}{なお}して}\R{求}{もと}めます。

\begin{bkfigure}
\begin{tikzpicture}[x=1mm, y=1mm]
  \node[draw=bkTheme, line width=1pt, dash pattern=on 2.4pt off 1.6pt, fill=white, rounded corners=2mm, minimum width=40mm, minimum height=27mm] (n1) at (21,0) {};
  \node[draw=bkTheme, line width=1pt, fill=bkTheme!8, rounded corners=2mm, minimum width=40mm, minimum height=27mm] (n2) at (70,0) {};
  \node[draw=bkTheme, line width=1pt, fill=bkTheme, rounded corners=2mm, minimum width=40mm, minimum height=27mm] (n3) at (119,0) {};
  \node[anchor=north, font=\sffamily\bfseries\small, text=bkInk] at ([yshift=-1.5mm]n1.north) {①\,\R{名前}{な|まえ}の\R{段}{だん}};
  \node[anchor=north, font=\sffamily\bfseries\small, text=bkInk] at ([yshift=-1.5mm]n2.north) {②\,\R{直}{なお}す\R{段}{だん}};
  \node[anchor=north, font=\sffamily\bfseries\small, text=white] at ([yshift=-1.5mm]n3.north) {③\,\R{楽}{らく}をする\R{段}{だん}};
  \node[anchor=north, text width=36mm, align=left, font=\sffamily\footnotesize, text=bkInk] at ([yshift=-8.5mm]n1.north) {\R{記号}{き|ごう}を\R{読}{よ}む：\R{何}{なに}を・どこで・1\R{枚}{まい}の\R{量}{りょう}は\R{何}{なに}か};
  \node[anchor=north, text width=36mm, align=left, font=\sffamily\footnotesize, text=bkInk] at ([yshift=-8.5mm]n2.north) {\R{場所}{ば|しょ}を\R{数}{かず}で\R{表}{あらわ}し、\R{範囲}{はん|い}を\R{決}{き}めて、ふつうの\R{計算}{けい|さん}に\R{直}{なお}す};
  \node[anchor=north, text width=36mm, align=left, font=\sffamily\footnotesize, text=white] at ([yshift=-8.5mm]n3.north) {\R{直}{なお}す\R{前}{まえ}に\R{確}{たし}かめる：\R{対称}{たい|しょう}・\R{分}{わ}ける・\R{定理}{てい|り}・\R{数値}{すう|ち}};
  \draw[-{Stealth[length=2.2mm]}, line width=1.1pt, bkGray] ([xshift=0.8mm]n1.east) -- ([xshift=-0.8mm]n2.west);
  \draw[-{Stealth[length=2.2mm]}, line width=1.1pt, bkGray, dashed] ([xshift=-0.8mm]n3.west) -- ([xshift=0.8mm]n2.east);
  \node[anchor=north, font=\sffamily\footnotesize, text=bkInk, inner sep=0.5pt, align=flush center] at ($(n2.east)!0.5!(n3.west)+(0,-4.5)$) {\R{直}{なお}す\\\R{前}{まえ}に\R{見}{み}る};
\end{tikzpicture}
\bkfigcaption{3つの\R{段}{だん}の\R{札}{ふだ}。\R{点線}{てん|せん}の\R{枠}{わく}・\R{実線}{じっ|せん}の\R{枠}{わく}・ぬりつぶしの\R{形}{かたち}でも\R{見分}{み|わ}けられる}
\end{bkfigure}

\R{第}{だい}8\R{巻}{かん}の\R{例}{れい}で\R{見}{み}てみましょう。\R{横}{よこ} 2、たて 1 の\R{長方形}{ちょう|ほう|けい}の\R{板}{いた}を、$x$ が $0$ から $2$、$y$ が $0$ から $1$ の\R{所}{ところ}に\R{置}{お}きます。$\sigma$ は\R{板}{いた}の1\,cm²あたりの\R{重}{おも}さのような\R{量}{りょう}です。

\begin{bkcenter}
\begin{bktable}{colspec={Q[l,m,wd=27mm] Q[l,m,wd=97mm]}, rowsep=2pt, cells={font=\small}, row{1}={font=\sffamily\bfseries\small}}
\R{段}{だん} & \R{書}{か}くこと \\
①\,\R{名前}{な|まえ}の\R{段}{だん} & $\displaystyle\iint_D\sigma\,\dd A$：\R{板}{いた}の\R{上}{うえ}で「$\sigma$×タイル」をたす。 \\
②\,\R{直}{なお}す\R{段}{だん} & \R{内側}{うち|がわ}は $y$ を\R{止}{と}めて $x$ で $0$ から $2$ まで、\R{外側}{そと|がわ}は $y$ で $0$ から $1$ まで。→ ふつうの\R{定積分}{てい|せき|ぶん}（\R{第}{だい}5\R{巻}{かん}）を2\R{回}{かい}\R{計算}{けい|さん}する。 \\
\end{bktable}
\bktabcaption{「\R{名前}{な|まえ}」を「ふつうの\R{定積分}{てい|せき|ぶん}」に\R{直}{なお}す（\R{第}{だい}8\R{巻}{かん}の\R{書}{か}き\R{方}{かた}）}
\end{bkcenter}

\R{直}{なお}す\R{先}{さき}の\R{計算}{けい|さん}は、\R{第}{だい}8\R{巻}{かん}では、どれも「ふつうの\R{定積分}{てい|せき|ぶん}（\R{第}{だい}5\R{巻}{かん}）のくり\R{返}{かえ}し」です。\R{第}{だい}9\R{巻}{かん}では、ふつうの\R{定積分}{てい|せき|ぶん}・\R{表}{ひょう}を\R{引}{ひ}く・\R{部分分数}{ぶ|ぶん|ぶん|すう}・\R{行列}{ぎょう|れつ}の\R{計算}{けい|さん}の4つの\R{道具}{どう|ぐ}のどれかに\R{直}{なお}します。

どの\R{道具}{どう|ぐ}に\R{直}{なお}せばよいかは、\R{場合}{ば|あい}によってちがいます。そこで\R{第}{だい}8・9\R{巻}{かん}には、\textbf{\R{場合分}{ば|あい|わ}けの\R{表}{ひょう}}と\textbf{\R{流}{なが}れ\R{図}{ず}}があります。\R{流}{なが}れ\R{図}{ず}は、ひし\R{形}{がた}の「はい・いいえ」の\R{問}{と}いと、\R{四角}{し|かく}の\R{手順}{て|じゅん}だけでできた\R{図}{ず}です（この\R{本}{ほん}の7\R{節}{せつ}にもあります）。

\R{場合分}{ば|あい|わ}けの\R{表}{ひょう}と\R{流}{なが}れ\R{図}{ず}は\R{見開}{み|ひら}きにならべてあり、その\R{章}{しょう}の\R{例題}{れい|だい}は、どれも\R{流}{なが}れ\R{図}{ず}でたどれるようにしてあります。\R{第}{だい}8\R{巻}{かん}の\R{巻末}{かん|まつ}には「\R{積分}{せき|ぶん}の\R{地図}{ち|ず}」、\R{第}{だい}9\R{巻}{かん}の\R{巻末}{かん|まつ}には「\R{変換}{へん|かん}の\R{地図}{ち|ず}」があり、\R{名前}{な|まえ}と\R{実際}{じっ|さい}の\R{操作}{そう|さ}を1か\R{所}{しょ}で\R{見}{み}わたせます。

\begin{bkcenter}
\begin{bktable}{colspec={Q[c,m,wd=12mm] Q[l,m,wd=52mm] Q[l,m,wd=58mm]}, rowsep=1.6pt, cells={font=\small}, row{1}={font=\sffamily\bfseries\small}}
\R{巻}{かん} & 3つの\R{段}{だん}（\R{目}{め}・\R{手}{て}・\R{式}{しき}） & \R{名前}{な|まえ}と\R{実際}{じっ|さい}の\R{操作}{そう|さ} \\
\kgavol{6} & 6-10〜6-14 & ― \\
\kgavol{7} & 7-7〜7-10 & ― \\
\kgavol{8} & 8-1〜8-13（8-5 は2つの\R{段}{だん}） & 8-0（\R{予告}{よ|こく}）、8-6〜8-11、8-14、\R{巻末}{かん|まつ}の「\R{積分}{せき|ぶん}の\R{地図}{ち|ず}」 \\
\kgavol{9} & 9-1〜9-11 & すべての\R{章}{しょう}と、\R{巻末}{かん|まつ}の「\R{変換}{へん|かん}の\R{地図}{ち|ず}」 \\
\end{bktable}
\bktabcaption{2つの\R{型}{かた}を\R{使}{つか}っている\R{章}{しょう}（\R{第}{だい}6・7\R{巻}{かん}は3つの\R{段}{だん}だけ）}
\end{bkcenter}

「\R{名前}{な|まえ}の\R{段}{だん}・\R{直}{なお}す\R{段}{だん}・\R{楽}{らく}をする\R{段}{だん}」は、「\R{目}{め}・\R{手}{て}・\R{式}{しき}」の3つの\R{段}{だん}とは\R{別}{べつ}のものです。\R{目}{め}・\R{手}{て}・\R{式}{しき}のどの\R{段}{だん}の\R{中}{なか}でも、\R{計算}{けい|さん}をするときに\R{使}{つか}う、もう1つ\R{小}{ちい}さな\R{手順}{て|じゅん}です。\R{第}{だい}9\R{巻}{かん}の\R{終章}{しゅう|しょう}（9-13）で、\R{全}{ぜん}9\R{巻}{かん}の「\R{名前}{な|まえ}」と「\R{実際}{じっ|さい}の\R{操作}{そう|さ}」をまとめます。

\bkfukidashi{delta}{\R{記号}{き|ごう}を\R{見}{み}たら、まず「これは\R{何}{なに}の\R{名前}{な|まえ}？」と\R{考}{かんが}えよう。\R{次}{つぎ}に「\R{知}{し}っている\R{計算}{けい|さん}のどれに\R{直}{なお}せる？」だよ。}


%% =====================================================================
%% 6　案内役のソラ・ミナ・デルタ（2ページ）
%% =====================================================================
\kgasec{6}{\R{案内役}{あん|ない|やく}のソラ・ミナ・デルタ}
  {シリーズを\R{通}{とお}して、3\R{人}{にん}がいっしょに\R{歩}{ある}きます。\R{困}{こま}ったときは、3\R{人}{にん}のことばをヒントにしよう。}
  {案内役のソラ・ミナ・デルタ}

\begin{bkfigure}
\begin{tikzpicture}[x=1mm, y=1mm]
  \foreach \x in {22,70,118} {\fill[bkTheme!7, rounded corners=3mm] (\x-22,0) rectangle (\x+22,96);}
  \begin{scope}\clip[rounded corners=3mm] (0,0) rectangle (44,96);\bkChar{sora}{(22,40)}{56mm}\end{scope}
  \begin{scope}\clip[rounded corners=3mm] (48,0) rectangle (92,96);\bkChar{mina}{(70,40)}{56mm}\end{scope}
  \bkChar{delta}{(118,48)}{34mm}
  \foreach \x/\n in {22/ソラ, 70/ミナ, 118/デルタ}
    \node[fill=bkTheme, text=white, rounded corners=2mm, inner xsep=4mm, inner ysep=1.4mm, font=\sffamily\bfseries\small] at (\x,39) {\n};
  \foreach \x/\t in {22/{「なぜ？」と\R{聞}{き}くのが\R{大好}{だい|す}き。たしかめるまで、なっとくしない。},
                     70/{\R{図}{ず}やグラフをかくのがとくい。\R{手}{て}を\R{動}{うご}かして\R{考}{かんが}える。},
                     118/{\R{計算}{けい|さん}と\R{記録}{き|ろく}がとくいなロボット。\R{名前}{な|まえ}は「\R{変化}{へん|か}」を\R{表}{あらわ}す\R{記号}{き|ごう}Δから。}}
    \node[anchor=north, text width=38mm, align=left, font=\small, text=bkInk] at (\x,33) {\t};
\end{tikzpicture}
\bkfigcaption{\R{案内役}{あん|ない|やく}の3\R{人}{にん}（\R{各巻}{かく|かん}の\R{表紙}{ひょう|し}の\R{裏}{うら}でも\R{紹介}{しょう|かい}しています）}
\end{bkfigure}

\kgaitem{ソラとミナは、\R{巻}{かん}が\R{進}{すす}むごとに\R{成長}{せい|ちょう}します。\R{第}{だい}7\R{巻}{かん}からは\R{大学生}{だい|がく|せい}です。}
\kgaitem{デルタの\R{名前}{な|まえ}の Δ（デルタ）は、\R{変}{か}わった\R{量}{りょう}を\R{表}{あらわ}す\R{記号}{き|ごう}です。\R{第}{だい}2\R{巻}{かん}の 2-3 で\R{習}{なら}います。}
\kgaitem{3\R{人}{にん}は、\R{表紙}{ひょう|し}の\R{裏}{うら}・\R{序章}{じょ|しょう}・\R{部}{ぶ}の\R{扉}{とびら}、そして\R{章}{しょう}の\R{中}{なか}の「ふきだし」に\R{出}{で}てきます。}

\clearpage
\bksubsection{1}{ふきだし}

\R{章}{しょう}の\R{中}{なか}では、3\R{人}{にん}は\R{次}{つぎ}のような「ふきだし」で\R{話}{はな}します。\R{次}{つぎ}の\R{章}{しょう}へのつながりや、\R{考}{かんが}えるヒントを\R{教}{おし}えてくれます。しゃべりすぎないように、ふきだしは\R{見開}{み|ひら}き（2ページ）に1\R{回}{かい}までにしています。

\bkfukidashi{mina}{「すぐ」「たくさん」「とても\R{重}{おも}い」も、\R{数}{かず}と\R{単位}{たん|い}で\R{言}{い}いかえられるかな？\quad この\R{巻}{かん}で、そのことばを1つずつ\R{学}{まな}ぼう。}
\vspace{-2mm}
{\centering\sffamily\small{\color{\bkTT}\bfseries \R{図}{ず}\refstepcounter{bkfig}\thebkfig}\hspace{0.8\zw}ふきだしの\R{例}{れい}（\R{第}{だい}2\R{巻}{かん}の\R{序章}{じょ|しょう} 2-0 から）\par}\vspace{3mm}

\bksubsection{2}{3\R{人}{にん}の\R{会話}{かい|わ}}

\R{序章}{じょ|しょう}では、3\R{人}{にん}の\R{会話}{かい|わ}で、その\R{巻}{かん}の\R{読}{よ}み\R{方}{かた}を\R{紹介}{しょう|かい}することがあります。\R{次}{つぎ}は、\R{第}{だい}2\R{巻}{かん}の\R{序章}{じょ|しょう}の\R{会話}{かい|わ}です。

\begin{bkfigure}
\begin{tikzpicture}[x=1mm, y=1mm]
  \kgatalk{0}{sora}{ソラ}{\R{第}{だい}1\R{巻}{かん}では、\R{数}{かぞ}えたり、はかったりしたね。\R{第}{だい}2\R{巻}{かん}では\R{何}{なに}をするの？}
  \kgatalk{-19}{mina}{ミナ}{0より\R{小}{ちい}さい\R{数}{かず}や、\R{文字}{も|じ}の $x$ も\R{出}{で}てくるんだって。むずかしそう……。}
  \kgatalk{-38}{delta}{デルタ}{だいじょうぶ。どの\R{章}{しょう}も、\R{場面}{ば|めん}を\R{見}{み}て、\R{手}{て}を\R{動}{うご}かして、\R{気}{き}づいてから\R{式}{しき}にするよ。わからなくなったら、\R{前}{まえ}の\R{章}{しょう}にもどればいいよ。}
\end{tikzpicture}
\bkfigcaption{\R{第}{だい}2\R{巻}{かん}の\R{序章}{じょ|しょう}（2-0）の\R{会話}{かい|わ}}
\end{bkfigure}

\bksubsection{3}{\R{部}{ぶ}の\R{扉}{とびら}}

\R{各巻}{かく|かん}の\R{部}{ぶ}の\R{始}{はじ}まりには、\R{巻}{かん}の\R{色}{いろ}の\R{扉}{とびら}があります。\R{扉}{とびら}には、その\R{部}{ぶ}で\R{答}{こた}える3つの「はてな？」と、デルタのひとことがのっています。\R{右}{みぎ}のページは「この\R{部}{ぶ}で\R{学}{まな}ぶこと」で、\R{部}{ぶ}の\R{章}{しょう}の\R{一覧}{いち|らん}です。

%% =====================================================================
%% 7　どこから読む？（4ページ）
%% =====================================================================
\kgasec{7}{どこから\R{読}{よ}む？}
  {いま\R{習}{なら}っていること・やりたいことから、はじめる\R{巻}{かん}を\R{選}{えら}ぼう。\R{迷}{まよ}ったら、1つ\R{前}{まえ}の\R{巻}{かん}から。}
  {どこから読む？}

このシリーズは、\R{第}{だい}1\R{巻}{かん}から\R{順}{じゅん}に\R{読}{よ}むと、\R{考}{かんが}えがいちばんよくつながるように\R{作}{つく}ってあります。\R{初}{はじ}めて\R{読}{よ}む\R{人}{ひと}や\R{迷}{まよ}った\R{人}{ひと}は、\R{第}{だい}1\R{巻}{かん}から\R{始}{はじ}めましょう。もう\R{知}{し}っていることが\R{多}{おお}い\R{人}{ひと}は、とちゅうの\R{巻}{かん}から\R{始}{はじ}めてもかまいません。

\bksubsection{1}{3\R{通}{とお}りの\R{始}{はじ}め\R{方}{かた}}

\begin{bkfigure}
\begin{tikzpicture}[x=1mm, y=1mm]
  \foreach \t/\s/\c [count=\i from 0] in {{\R{順番}{じゅん|ばん}に\R{読}{よ}む}/{\R{第}{だい}1\R{巻}{かん}から、\\\R{順}{じゅん}に\R{読}{よ}み\R{進}{すす}める}/1,
      {\R{習}{なら}っていることから}/{\R{流}{なが}れ\R{図}{ず}（\bkpageno{\kgaPage{kga-flowp}}）で、\\はじめる\R{巻}{かん}を\R{決}{き}める}/2,
      {やりたいことから}/{\R{表}{ひょう}（\bkpageno{\kgaPage{kga-wantp}}）で\R{章}{しょう}を\\\R{見}{み}つけ、\R{道}{みち}しるべで\R{前}{まえ}へ}/3} {
    \node[draw=bkTheme, line width=1pt, fill=white, rounded corners=2.5mm, minimum width=42mm, minimum height=29mm, anchor=north west] (w\i) at (\i*48,0) {};
    \node[anchor=north, circle, fill=bkTheme, text=white, font=\sffamily\bfseries\small, minimum size=6mm, inner sep=0pt] (wc\i) at ([yshift=-2mm]w\i.north) {\c};
    \node[anchor=north, font=\sffamily\bfseries\small, text=\bkTT, inner sep=0pt] (wt\i) at ([yshift=-1.6mm]wc\i.south) {\t};
    \node[anchor=north, text width=39mm, align=flush center, font=\sffamily\footnotesize, text=bkInk, inner sep=0pt] at ([yshift=-1.8mm]wt\i.south) {\s};}
\end{tikzpicture}
\bkfigcaption{3\R{通}{とお}りの\R{始}{はじ}め\R{方}{かた}}
\end{bkfigure}


どの\R{始}{はじ}め\R{方}{かた}でも、\R{章}{しょう}のはじめの\R{道}{みち}しるべと、\R{思}{おも}い\R{出}{だ}そうが\R{助}{たす}けてくれます。\R{前}{まえ}の\R{巻}{かん}で\R{学}{まな}ぶことは、\R{思}{おも}い\R{出}{だ}そうで\R{短}{みじか}くふり\R{返}{かえ}るだけなので、わからないときは\R{札}{ふだ}の\R{章}{しょう}にもどって\R{読}{よ}みましょう。


\bksubsection{2}{とちゅうの\R{巻}{かん}から\R{始}{はじ}めるときのコツ}

\kgaitem{まず、その\R{巻}{かん}のもくじ・\R{路線図}{ろ|せん|ず}・\R{序章}{じょ|しょう}を\R{見}{み}よう。\R{序章}{じょ|しょう}には、その\R{巻}{かん}で\R{学}{まな}ぶことと\R{読}{よ}み\R{方}{かた}がまとめてあります。}
\kgaitem{わからないことばは、\R{各巻}{かく|かん}の\R{巻末}{かん|まつ}の「さくいん」で\R{引}{ひ}けます。\R{第}{だい}9\R{巻}{かん}には、\R{第}{だい}1〜9\R{巻}{かん}をまとめた「シリーズ\R{全体}{ぜん|たい}のさくいん」もあります。}
\kgaitem{\R{新}{あたら}しい\R{記号}{き|ごう}は、\R{各巻}{かく|かん}の\R{巻末}{かん|まつ}の「\R{記号}{き|ごう}カード」の\R{一覧}{いち|らん}で、\R{読}{よ}み\R{方}{かた}と\R{初}{はじ}めて\R{出}{で}た\R{章}{しょう}がわかります。}



\clearpage
\bksubsection{3}{やりたいことから\R{選}{えら}ぶ}\label{kga-wantp}

「ロボットを\R{動}{うご}かしたい」「\R{音}{おと}を\R{調}{しら}べたい」のように、やりたいことが\R{決}{き}まっている\R{人}{ひと}は、\R{下}{した}の\R{表}{ひょう}から\R{章}{しょう}を\R{選}{えら}んでもかまいません。\R{章}{しょう}の\R{内容}{ない|よう}は、それぞれの「\R{物理}{ぶつ|り}・\R{工学}{こう|がく}ではこう\R{使}{つか}う」などで\R{取}{と}り\R{上}{あ}げている\R{話題}{わ|だい}です。

\begin{bkcenter}
\begin{bktable}{colspec={Q[l,m,wd=36mm] Q[l,m,wd=88mm]}, rowsep=1.6pt, cells={font=\small}, row{1}={font=\sffamily\bfseries\small}}
やりたいこと & \R{章}{しょう}（\R{章}{しょう}の\R{題}{だい}・\R{話題}{わ|だい}） \\
ロボットの\R{腕}{うで}を\R{動}{うご}かしたい & \kgapill{4-0} ロボットの\R{手}{て}は、どこにある？\quad\kgapill{4-7} \R{正弦定理}{せい|げん|てい|り}と\R{余弦定理}{よ|げん|てい|り}（\R{順運動学}{じゅん|うん|どう|がく}と\R{逆運動学}{ぎゃく|うん|どう|がく}）\quad\kgapill{7-6} \R{回転行列}{かい|てん|ぎょう|れつ}\quad\kgapill{7-11} \R{行列}{ぎょう|れつ}の\R{使}{つか}いみち \\
\R{電気}{でん|き}の\R{回路}{かい|ろ}を\R{知}{し}りたい & \kgapill{2-15} オームの\R{法則}{ほう|そく}\quad\kgapill{6-4} コンデンサーとコイルの\R{回路}{かい|ろ}\quad\kgapill{6-10} \R{交流}{こう|りゅう}の\R{計算}{けい|さん}\quad\kgapill{6-13} \R{共振}{きょう|しん} \\
\R{音}{おと}や\R{波}{なみ}を\R{調}{しら}べたい & \kgapill{4-9} うなり\quad\kgapill{9-0} \R{音}{おと}を\R{分}{わ}ける\quad\kgapill{9-1} フーリエ\R{級数}{きゅう|すう}\quad\kgapill{9-4} フーリエ\R{変換}{へん|かん}と FFT \\
\R{機械}{き|かい}の\R{動}{うご}きを\R{予想}{よ|そう}・\R{制御}{せい|ぎょ}したい & \kgapill{6-0} \R{未来}{み|らい}を\R{予想}{よ|そう}する\quad\kgapill{6-7} \R{減衰振動}{げん|すい|しん|どう}\quad\kgapill{9-8} \R{伝達関数}{でん|たつ|かん|すう}\quad\kgapill{9-9} \R{極}{きょく}と\R{安定性}{あん|てい|せい}\quad\kgapill{9-10} \R{状態方程式}{じょう|たい|ほう|てい|しき} \\
\R{熱}{ねつ}の\R{伝}{つた}わり\R{方}{かた}を\R{知}{し}りたい & \kgapill{6-4} \R{冷}{さ}めるお\R{湯}{ゆ}\quad\kgapill{8-12} ラプラシアン\quad\kgapill{8-13} \R{熱}{ねつ}と\R{波}{なみ}の\R{伝}{つた}わり\R{方}{かた} \\
\R{電気}{でん|き}と\R{磁気}{じ|き}を\R{知}{し}りたい & \kgapill{4-14} \R{外積}{がい|せき}（\R{電気}{でん|き}モーター）\quad\kgapill{8-9} \R{発散}{はっ|さん}\quad\kgapill{8-10} \R{回転}{かい|てん}\quad\kgapill{8-11} ガウスの\R{定理}{てい|り}\quad\kgapill{8-14} \R{電気}{でん|き}と\R{磁気}{じ|き}の4つの\R{式}{しき} \\
\R{空気}{くう|き}や\R{水}{みず}の\R{流}{なが}れを\R{知}{し}りたい & \kgapill{8-5} ベクトル\R{場}{ば}\quad\kgapill{8-9} \R{連続}{れん|ぞく}の\R{式}{しき}\quad\kgapill{8-14} ナビエ・ストークス\R{方程式}{ほう|てい|しき} \\
はかったデータを\R{調}{しら}べたい & \kgapill{2-6} はかった\R{数}{かず}のずれ\quad\kgapill{3-13} ヒストグラム\quad\kgapill{3-14} \R{標準偏差}{ひょう|じゅん|へん|さ}\quad\kgapill{3-16} \R{最小二乗法}{さい|しょう|に|じょう|ほう} \\
コンピュータで\R{計算}{けい|さん}したい & \kgapill{2-17} 2\R{進数}{しん|すう}\quad\kgapill{2-18} \R{表計算}{ひょう|けい|さん}\quad\kgapill{2-19} プログラム\quad\kgapill{6-2} オイラー\R{法}{ほう} \\
\end{bktable}
\bktabcaption{やりたいことと\R{章}{しょう}}
\end{bkcenter}

\R{選}{えら}んだ\R{章}{しょう}の★が\R{高}{たか}くて\R{読}{よ}めないときは、\R{章}{しょう}のはじめの\R{道}{みち}しるべで、わかるところまで\R{前}{まえ}の\R{章}{しょう}にさかのぼり、そこから\R{順番}{じゅん|ばん}にもどって\R{読}{よ}もう。

\clearpage
\bksubsection{4}{\R{流}{なが}れ\R{図}{ず}ではじめる\R{巻}{かん}を\R{決}{き}める}\label{kga-flowp}

「スタート」から、ひし\R{形}{がた}の\R{問}{と}いに「はい」か「いいえ」で\R{答}{こた}えて\R{進}{すす}みます。「いいえ」の\R{先}{さき}が、はじめる\R{巻}{かん}と\R{章}{しょう}です。「\R{習}{なら}い\R{始}{はじ}めた？」は、いま\R{習}{なら}っている\R{所}{ところ}でも「はい」です。\R{大人}{おと|な}の\R{人}{ひと}は、\R{学}{まな}び\R{直}{なお}したい\R{所}{ところ}を「いま\R{習}{なら}っている」と\R{考}{かんが}えて\R{答}{こた}えましょう。\R{迷}{まよ}ったら、1つ\R{前}{まえ}の\R{巻}{かん}から\R{始}{はじ}めよう。

\begin{bkfigure}
\begin{tikzpicture}[x=1mm, y=1mm]
  \node[kga box, minimum width=24mm] (s) at (51,0) {スタート};
  \node[kga q1] (q1) at (51,-19) {\R{九九}{く|く}を、\R{習}{なら}い\R{始}{はじ}めた？};
  \node[kga go=1, anchor=west, align=flush center, minimum width=22mm, inner ysep=0.8mm] (g1) at (114,-19) {\R{第}{だい}1\R{巻}{かん}\\1-0 から};
  \draw[kga arr] (q1.east) -- node[kga yn, above, yshift=0.5mm] {いいえ} (g1.west);
  \node[kga q1] (q2) at (51,-41.5) {\R{分度器}{ぶん|ど|き}で\R{角}{かく}をはかることを、\R{習}{なら}い\R{始}{はじ}めた？};
  \node[kga go=1, anchor=west, align=flush center, minimum width=22mm, inner ysep=0.8mm] (g2) at (114,-41.5) {\R{第}{だい}1\R{巻}{かん}\\1-2 から};
  \draw[kga arr] (q2.east) -- node[kga yn, above, yshift=0.5mm] {いいえ} (g2.west);
  \node[kga q1] (q3) at (51,-64) {\R{小数}{しょう|すう}のかけ\R{算}{ざん}を、\R{習}{なら}い\R{始}{はじ}めた？};
  \node[kga go=1, anchor=west, align=flush center, minimum width=22mm, inner ysep=0.8mm] (g3) at (114,-64) {\R{第}{だい}1\R{巻}{かん}\\1-15 から};
  \draw[kga arr] (q3.east) -- node[kga yn, above, yshift=0.5mm] {いいえ} (g3.west);
  \node[kga q1] (q4) at (51,-86.5) {\R{分数}{ぶん|すう}のかけ\R{算}{ざん}を、\R{習}{なら}い\R{始}{はじ}めた？};
  \node[kga go=2, anchor=west, align=flush center, minimum width=22mm, inner ysep=0.8mm] (g4) at (114,-86.5) {\R{第}{だい}2\R{巻}{かん}\\2-0 から};
  \draw[kga arr] (q4.east) -- node[kga yn, above, yshift=0.5mm] {いいえ} (g4.west);
  \node[kga q1] (q5) at (51,-109) {マイナスの\R{数}{かず}の\R{計算}{けい|さん}を、\R{習}{なら}い\R{始}{はじ}めた？};
  \node[kga go=2, anchor=west, align=flush center, minimum width=22mm, inner ysep=0.8mm] (g5) at (114,-109) {\R{第}{だい}2\R{巻}{かん}\\2-4 から};
  \draw[kga arr] (q5.east) -- node[kga yn, above, yshift=0.5mm] {いいえ} (g5.west);
  \node[kga q1] (q6) at (51,-131.5) {\R{連立方程式}{れん|りつ|ほう|てい|しき}を、\R{習}{なら}い\R{始}{はじ}めた？};
  \node[kga go=2, anchor=west, align=flush center, minimum width=22mm, inner ysep=0.8mm] (g6) at (114,-131.5) {\R{第}{だい}2\R{巻}{かん}\\2-2 から};
  \draw[kga arr] (q6.east) -- node[kga yn, above, yshift=0.5mm] {いいえ} (g6.west);
  \node[kga box, draw=bkTheme, minimum width=34mm] (n) at (51,-152) {\R{右}{みぎ}のページの \textbf{A} へ};
  \draw[kga arr] (s) -- (q1);
  \draw[kga arr] (q1.south) -- node[kga yn, right, xshift=1mm] {はい} (q2.north);
  \draw[kga arr] (q2.south) -- node[kga yn, right, xshift=1mm] {はい} (q3.north);
  \draw[kga arr] (q3.south) -- node[kga yn, right, xshift=1mm] {はい} (q4.north);
  \draw[kga arr] (q4.south) -- node[kga yn, right, xshift=1mm] {はい} (q5.north);
  \draw[kga arr] (q5.south) -- node[kga yn, right, xshift=1mm] {はい} (q6.north);
  \draw[kga arr] (q6.south) -- node[kga yn, right, xshift=1mm] {はい} (n.north);
\end{tikzpicture}
\bkfigcaption{はじめる\R{巻}{かん}・\R{章}{しょう}を\R{決}{き}める\R{流}{なが}れ\R{図}{ず}（1）}
\end{bkfigure}

\clearpage
\begin{bkfigure}
\begin{tikzpicture}[x=1mm, y=1mm]
  \node[kga box, draw=bkTheme, minimum width=24mm] (s) at (51,0) {\textbf{A}};
  \node[kga q1] (q7) at (51,-19) {\R{平方根}{へい|ほう|こん}（$\sqrt{\phantom{x}}$）を、\R{習}{なら}い\R{始}{はじ}めた？};
  \node[kga go=2, anchor=west, align=flush center, minimum width=22mm, inner ysep=0.8mm] (g7) at (114,-19) {\R{第}{だい}2\R{巻}{かん}\\2-12 から};
  \draw[kga arr] (q7.east) -- node[kga yn, above, yshift=0.5mm] {いいえ} (g7.west);
  \node[kga q1] (q8) at (51,-41.5) {2\R{次関数}{じ|かん|すう}か\R{三角比}{さん|かく|ひ}を、\R{習}{なら}い\R{始}{はじ}めた？};
  \node[kga go=2, anchor=west, align=flush center, minimum width=22mm, inner ysep=0.8mm] (g8) at (114,-41.5) {\R{第}{だい}2\R{巻}{かん}\\2-6 から};
  \draw[kga arr] (q8.east) -- node[kga yn, above, yshift=0.5mm] {いいえ} (g8.west);
  \node[kga q1] (q9) at (51,-64) {\R{三角関数}{さん|かく|かん|すう}か\R{指数}{し|すう}・\R{対数}{たい|すう}を、\R{習}{なら}い\R{始}{はじ}めた？};
  \node[kga go=2, anchor=west, align=flush center, minimum width=22mm, inner ysep=0.8mm] (g9) at (114,-64) {\R{第}{だい}2\R{巻}{かん}\\2-9 から};
  \draw[kga arr] (q9.east) -- node[kga yn, above, yshift=0.5mm] {いいえ} (g9.west);
  \node[kga q1] (q10) at (51,-86.5) {\R{指数}{し|すう}と\R{対数}{たい|すう}を、\R{習}{なら}い\R{始}{はじ}めた？};
  \node[kga go=4, anchor=west, align=flush center, minimum width=22mm, inner ysep=0.8mm] (g10) at (114,-86.5) {\R{第}{だい}4\R{巻}{かん}\\4-2 から};
  \draw[kga arr] (q10.east) -- node[kga yn, above, yshift=0.5mm] {いいえ} (g10.west);
  \node[kga q1] (q11) at (51,-109) {\R{微分}{び|ぶん}と\R{積分}{せき|ぶん}を、\R{習}{なら}い\R{始}{はじ}めた？};
  \node[kga go=2, anchor=west, align=flush center, minimum width=22mm, inner ysep=0.8mm] (g11) at (114,-109) {\R{第}{だい}2\R{巻}{かん}\\2-7 から};
  \draw[kga arr] (q11.east) -- node[kga yn, above, yshift=0.5mm] {いいえ} (g11.west);
  \node[kga q1] (q12) at (51,-131.5) {$e^x$ の\R{微分}{び|ぶん}を、\R{習}{なら}い\R{始}{はじ}めた？};
  \node[kga go=5, anchor=west, align=flush center, minimum width=22mm, inner ysep=0.8mm] (g12) at (114,-131.5) {\R{第}{だい}5\R{巻}{かん}\\5-0 から};
  \draw[kga arr] (q12.east) -- node[kga yn, above, yshift=0.5mm] {いいえ} (g12.west);
  \node[kga q1] (q13) at (51,-154) {\R{大学}{だい|がく}の\R{数学}{すう|がく}を、\R{習}{なら}い\R{始}{はじ}めた？};
  \node[kga go=5, anchor=west, align=flush center, minimum width=22mm, inner ysep=0.8mm] (g13) at (114,-154) {\R{第}{だい}5\R{巻}{かん}\\5-3 から};
  \draw[kga arr] (q13.east) -- node[kga yn, above, yshift=0.5mm] {いいえ} (g13.west);
  \node[kga box, text width=118mm, align=left, anchor=north] (n) at (68,-174.5) {\R{大学}{だい|がく}の\R{数学}{すう|がく}を\R{習}{なら}っている\R{人}{ひと}や\R{大人}{おと|な}の\R{人}{ひと}は、\R{学}{まな}びたい\R{所}{ところ}の\R{巻}{かん}から（\R{微分}{び|ぶん}・\R{積分}{せき|ぶん}は\R{第}{だい}5\R{巻}{かん}、\R{微分方程式}{び|ぶん|ほう|てい|しき}は\R{第}{だい}6\R{巻}{かん}、\R{行列}{ぎょう|れつ}は\R{第}{だい}7\R{巻}{かん}）。\R{工学}{こう|がく}の\R{専門書}{せん|もん|しょ}を\R{読}{よ}みたい\R{人}{ひと}は\R{第}{だい}6\R{巻}{かん}から。};
  \draw[kga arr] (s) -- (q7);
  \draw[kga arr] (q7.south) -- node[kga yn, right, xshift=1mm] {はい} (q8.north);
  \draw[kga arr] (q8.south) -- node[kga yn, right, xshift=1mm] {はい} (q9.north);
  \draw[kga arr] (q9.south) -- node[kga yn, right, xshift=1mm] {はい} (q10.north);
  \draw[kga arr] (q10.south) -- node[kga yn, right, xshift=1mm] {はい} (q11.north);
  \draw[kga arr] (q11.south) -- node[kga yn, right, xshift=1mm] {はい} (q12.north);
  \draw[kga arr] (q12.south) -- node[kga yn, right, xshift=1mm] {はい} (q13.north);
  \draw[kga arr] (q13.south) -- node[kga yn, right, xshift=1mm] {はい} (n.north);
\end{tikzpicture}
\bkfigcaption{はじめる\R{巻}{かん}・\R{章}{しょう}を\R{決}{き}める\R{流}{なが}れ\R{図}{ず}（2）}
\end{bkfigure}



%% =====================================================================
%% 8　勉強のしかた（2ページ）
%% =====================================================================
\kgasec{8}{\R{勉強}{べん|きょう}のしかた}
  {\R{電卓}{でん|たく}・\R{表計算}{ひょう|けい|さん}・プログラムを\R{使}{つか}う\R{時期}{じ|き}、こたえの\R{見}{み}\R{方}{かた}、こまったときにすること。1\R{回}{かい}でわからなくても、だいじょうぶ。}
  {勉強のしかた}

\bksubsection{1}{\R{使}{つか}う\R{道具}{どう|ぐ}}

\begin{bkfigure}
\begin{tikzpicture}[x=1mm, y=1mm]
  \foreach \v in {1,...,9} {\node[inner sep=0pt] at ({34+(\v-0.5)*11.5},44) {\kgavol{\v}};
    \draw[bkLine, line width=0.4pt] ({34+(\v-1)*11.5},-1) -- ({34+(\v-1)*11.5},40);}
  \draw[bkLine, line width=0.4pt] (137.5,-1) -- (137.5,40);
  \foreach \t/\a/\b/\c/\l [count=\i from 0] in {
      {\R{本}{ほん}に\R{書}{か}きこむ}/0/1/bkOrange/{},
      {ノートに\R{解}{と}く}/1/9/bkBlue/{},
      {\R{電卓}{でん|たく}}/1.05/9/bkGreen/{2-1から},
      {\R{表計算}{ひょう|けい|さん}}/1.9/9/bkPurple/{2-18から},
      {プログラム}/1.95/9/bkNavy/{2-19から}} {
    \pgfmathsetmacro\kgay{34-\i*8}
    \node[anchor=east, font=\sffamily\footnotesize, text=bkInk, inner sep=0pt] at (32,\kgay) {\t};
    \fill[\c, rounded corners=1.2mm] ({34+\a*11.5+0.8},\kgay-2.6) rectangle ({34+\b*11.5-0.8},\kgay+2.6);
    \node[anchor=west, font=\sffamily\footnotesize, text=white, inner sep=0pt] at ({34+\a*11.5+2},\kgay) {\l};}
\end{tikzpicture}
\bkfigcaption{\R{巻}{かん}ごとに\R{使}{つか}う\R{道具}{どう|ぐ}（\R{帯}{おび}のある\R{巻}{かん}で\R{使}{つか}う）}
\end{bkfigure}

\kgaitem{\R{第}{だい}1\R{巻}{かん}は、\R{本}{ほん}に\R{書}{か}きこんで\R{使}{つか}います。\R{数字}{すう|じ}を\R{書}{か}く\R{練習}{れん|しゅう}もかねています。\R{第}{だい}2\R{巻}{かん}からは、ノートを1\R{冊}{さつ}\R{用意}{よう|い}して、\R{問題}{もん|だい}はノートに\R{解}{と}きます。\R{本}{ほん}に\R{書}{か}きこまないので、\R{何度}{なん|ど}でもやり\R{直}{なお}せます。}
\kgaitem{\R{電卓}{でん|たく}は 2-1 から\R{使}{つか}います。おす\R{前}{まえ}に\R{式}{しき}をノートに\R{書}{か}き、\R{答}{こた}えのだいたいの\R{大}{おお}きさが\R{合}{あ}っているかを\R{確}{たし}かめます。}
\kgaitem{\R{表計算}{ひょう|けい|さん}は 2-18 で\R{学}{まな}んでから\R{使}{つか}います。2-19 では、プログラム（ブロックをつないで\R{書}{か}く\R{方法}{ほう|ほう}と、Python などの\R{文字}{も|じ}で\R{書}{か}く\R{方法}{ほう|ほう}）を\R{学}{まな}びます。\R{第}{だい}6\R{巻}{かん}からは、\R{巻末}{かん|まつ}にプログラムのひな\R{形}{がた}がのっています。}
\kgaitem{コンピュータは、\R{手}{て}で\R{考}{かんが}えたあとの\R{確}{たし}かめと、くり\R{返}{かえ}しの\R{計算}{けい|さん}に\R{使}{つか}います。\R{途中}{と|ちゅう}の\R{式}{しき}を\R{書}{か}くのは、いつも\R{自分}{じ|ぶん}です。}

\bksubsection{2}{こたえの\R{見}{み}\R{方}{かた}}

\R{本}{ほん}の\R{後}{うし}ろの「こたえ」には、じぶんでとこうと\R{章末問題}{しょう|まつ|もん|だい}のこたえがのっています。こたえだけなので、\R{解}{と}き\R{方}{かた}は、\R{同}{おな}じ\R{形}{かたち}の\R{例題}{れい|だい}にもどって\R{確}{たし}かめます。くわしい\R{解}{と}き\R{方}{かた}は、これから\R{作}{つく}る\R{別}{べっ}\R{冊}{さつ}の\R{問題集}{もん|だい|しゅう}と\R{指導書}{し|どう|しょ}にのせます。

\kgaitem{まちがえたら、ノートの\R{式}{しき}を1\R{行}{ぎょう}ずつ\R{見}{み}なおして、どの\R{行}{ぎょう}でまちがえたかを\R{見}{み}つけよう。\R{単位}{たん|い}も、こたえの\R{一部}{いち|ぶ}です。}

\bksubsection{3}{こまったときは}

\begin{bkfigure}
\begin{tikzpicture}[x=1mm, y=1mm]
  \tikzset{kga qs/.style={kga qv, aspect=4.2, text width=34mm}, kga ans/.style={kga box, anchor=west, text width=43mm, align=left}}
  \node[kga box] (s) at (40,0) {\R{問題}{もん|だい}が\R{解}{と}けない};
  \node[kga qs] (q1) at (40,-17) {\R{同}{おな}じ\R{形}{かたち}の\R{例題}{れい|だい}を、\\\R{読}{よ}みなおした？};
  \node[kga qs] (q2) at (40,-41) {ふりかえりの\\「→ページ」を\R{読}{よ}んだ？};
  \node[kga qs] (q3) at (40,-65) {\R{道}{みち}しるべの\R{前}{まえ}の\R{章}{しょう}は、\\わかる？};
  \node[kga box, draw=bkTheme] (k) at (40,-85) {おうちの\R{人}{ひと}や\R{先生}{せん|せい}にきいてみよう};
  \node[kga qs] (q4) at (40,-105) {☆の\R{手前}{て|まえ}まで、\\\R{読}{よ}めた？};
  \node[kga ans] (a1) at (88,-17) {\R{例題}{れい|だい}の\R{式}{しき}を1\R{行}{ぎょう}ずつノートに\R{写}{うつ}してから、もう\R{一度}{いち|ど}とこう};
  \node[kga ans] (a2) at (88,-41) {そのページの\R{説明}{せつ|めい}と\R{図}{ず}を\R{読}{よ}んでから、もう\R{一度}{いち|ど}とこう};
  \node[kga ans] (a3) at (88,-65) {その\R{章}{しょう}にもどって\R{読}{よ}んでから、もう\R{一度}{いち|ど}とこう};
  \node[kga ans] (a5) at (88,-105) {いっしょに、\R{例題}{れい|だい}からもう\R{一度}{いち|ど}\R{読}{よ}もう};
  \node[kga box, text width=62mm, align=flush center] (e) at (40,-128) {\R{先}{さき}へ\R{進}{すす}もう。\\あとで、もう\R{一度}{いち|ど}もどってこよう};
  \draw[kga arr] (s) -- (q1);
  \draw[kga arr] (q1) -- node[kga yn, above, yshift=0.6mm] {いいえ} (a1);
  \draw[kga arr] (q1) -- node[kga yn, right, xshift=1mm] {はい} (q2);
  \draw[kga arr] (q2) -- node[kga yn, above, yshift=0.6mm] {いいえ} (a2);
  \draw[kga arr] (q2) -- node[kga yn, right, xshift=1mm] {はい} (q3);
  \draw[kga arr] (q3) -- node[kga yn, above, yshift=0.6mm] {いいえ} (a3);
  \draw[kga arr] (q3) -- node[kga yn, right, xshift=1mm] {はい} (k);
  \draw[kga arr] (k) -- (q4);
  \draw[kga arr] (q4) -- node[kga yn, above, yshift=0.6mm] {いいえ} (a5);
  \draw[kga arr] (q4) -- node[kga yn, right, xshift=1mm] {はい} (e);
\end{tikzpicture}
\bkfigcaption{こまったときの\R{流}{なが}れ\R{図}{ず}。1\R{回}{かい}でわからなくても、だいじょうぶ}
\end{bkfigure}

%% =====================================================================
%% 9　巻末の使い方と、その先へ（2ページ）
%% =====================================================================
\kgasec{9}{\R{巻末}{かん|まつ}の\R{使}{つか}い\R{方}{かた}と、その\R{先}{さき}へ}
  {\R{本}{ほん}の\R{後}{うし}ろには、\R{調}{しら}べるためのページがあります。\R{全}{ぜん}9\R{巻}{かん}を\R{終}{お}えたら、\R{物理}{ぶつ|り}シリーズへ。}
  {巻末の使い方と、その先へ}

\bksubsection{1}{\R{巻末}{かん|まつ}のページ}

\begin{bkcenter}
\begin{bktable}{colspec={Q[l,m,wd=56mm] Q[c,m,wd=5mm] Q[c,m,wd=5mm] Q[c,m,wd=5mm] Q[c,m,wd=5mm] Q[c,m,wd=5mm] Q[c,m,wd=5mm] Q[c,m,wd=5mm] Q[c,m,wd=5mm] Q[c,m,wd=5mm]}, colsep=1.2mm, rowsep=0pt, cells={font=\small}, row{1}={font=\sffamily\bfseries\small, abovesep=2pt}}
\R{巻末}{かん|まつ}のページ & \kgavol{1} & \kgavol{2} & \kgavol{3} & \kgavol{4} & \kgavol{5} & \kgavol{6} & \kgavol{7} & \kgavol{8} & \kgavol{9} \\
こたえ & ● & ● & ● & ● & ● & ● & ● & ● & ● \\
\R{記号}{き|ごう}カード（\R{一覧}{いち|らん}） & ● & ● & ● & ● & ● & ● & ● & ● & ● \\
よく\R{使}{つか}う\R{定数}{てい|すう}・\R{単位}{たん|い}・\R{公式}{こう|しき} & ▲ & ● & ● & ● & ● & ● & ● & ● & ● \\
\SetCell{l}\quad プログラムのひな\R{形}{がた} & & & & & & ● & ● & ● & ● \\
\SetCell{l}\quad \R{積分}{せき|ぶん}の\R{地図}{ち|ず}・\R{変換}{へん|かん}の\R{地図}{ち|ず} & & & & & & & & ● & ● \\
できるようになったことチェック\R{表}{ひょう} & ● & ● & ● & ● & ● & ● & ● & ● & ● \\
つぎの\R{巻}{かん}へ & ● & ● & ● & ● & ● & ● & ● & ● & ◆ \\
さくいん & ● & ● & ● & ● & ● & ● & ● & ● & ● \\
シリーズ\R{全体}{ぜん|たい}のさくいん & & & & & & & & & ● \\
ふろく & ● & & & & & & & & \\
\end{bktable}
\bktabcaption{\R{巻末}{かん|まつ}のページ（●＝ある。▲＝\R{第}{だい}1\R{巻}{かん}は「たんいのひょう」。◆＝\R{第}{だい}9\R{巻}{かん}は「\R{物理}{ぶつ|り}シリーズへ」。\R{第}{だい}1\R{巻}{かん}は「きごうカード」「チェックひょう」「つぎのステージへ」とひらがなの\R{名前}{な|まえ}）}
\end{bkcenter}

\kgaitem{\textbf{\R{記号}{き|ごう}カード}：\R{各章}{かく|しょう}の\R{記号}{き|ごう}カードを1か\R{所}{しょ}に\R{集}{あつ}めた\R{一覧}{いち|らん}です。\R{読}{よ}み\R{方}{かた}を\R{忘}{わす}れた\R{記号}{き|ごう}は、ここで\R{引}{ひ}けます。}
\kgaitem{\textbf{チェック\R{表}{ひょう}}：\R{各章}{かく|しょう}のふりかえりを1つの\R{表}{ひょう}にまとめたものです。できるようになったら、しるしをつけていこう。}
\kgaitem{\textbf{つぎの\R{巻}{かん}へ}：この\R{巻}{かん}の\R{道具}{どう|ぐ}が、この\R{先}{さき}の\R{巻}{かん}でどう\R{使}{つか}われるかの\R{予告}{よ|こく}です。}

\clearpage
\bksubsection{2}{\R{物理}{ぶつ|り}シリーズへ}

\R{全}{ぜん}9\R{巻}{かん}を\R{終}{お}えたら、\R{次}{つぎ}は\R{物理}{ぶつ|り}シリーズ（\R{全}{ぜん}5\R{巻}{かん}）です。\R{物理}{ぶつ|り}シリーズの\R{各章}{かく|しょう}の\R{道}{みち}しるべには、\R{前}{まえ}に\R{学}{まな}んでおく\R{数学}{すう|がく}の\R{章}{しょう}の\R{番号}{ばん|ごう}（6-2 など）が\R{書}{か}いてあります。\R{第}{だい}9\R{巻}{かん}の\R{巻末}{かん|まつ}「\R{物理}{ぶつ|り}シリーズへ」で、\R{物理}{ぶつ|り}の「なぜ？」とこのシリーズの\R{章}{しょう}の\R{対応}{たい|おう}を\R{見}{み}られます。

\begin{bkfigure}
\begin{tikzpicture}[x=1mm, y=1mm]
  \foreach \n/\t/\u [count=\i from 0] in {1/{\R{力学}{りき|がく}}/{8,7}, 2/{\R{熱学}{ねつ|がく}}/{8,3,5}, 3/{\R{波動}{は|どう}}/{9}, 4/{\R{電磁気学}{でん|じ|き|がく}}/{8,6}, 5/{\R{原子}{げん|し}}/{6,4,5}} {
    \node[draw=bkNavy, line width=1pt, fill=white, rounded corners=2.5mm, minimum width=25.5mm, minimum height=35mm, anchor=north west] (p\i) at (\i*28.6,0) {};
    \begin{scope}\clip[rounded corners=2.5mm] (p\i.north west) rectangle (p\i.south east);
      \fill[bkNavy] (p\i.north west) rectangle ([yshift=-14.5mm]p\i.north east);\end{scope}
    \node[anchor=north, text=white, font=\sffamily\bfseries\footnotesize, inner sep=0pt] at ([yshift=-2mm]p\i.north) {\R{物理}{ぶつ|り}\enspace\n};
    \node[anchor=north, text=white, font=\sffamily\bfseries\small, inner sep=0pt] at ([yshift=-8.2mm]p\i.north) {\t};
    \node[anchor=north, text=bkInk, font=\sffamily\footnotesize, inner sep=0pt] at ([yshift=-16.2mm]p\i.north) {よく\R{使}{つか}う};
    \node[anchor=north, text=bkInk, font=\sffamily\footnotesize, inner sep=0pt] at ([yshift=-20.7mm]p\i.north) {\R{数学}{すう|がく}の\R{巻}{かん}};
    \node[anchor=north, align=flush center, inner sep=0pt] at ([yshift=-26.6mm]p\i.north) {\foreach \w [count=\j] in \u {\ifnum\j>1\hspace{0.6mm}\fi\kgavol{\w}}};}
\end{tikzpicture}
\bkfigcaption{\R{物理}{ぶつ|り}シリーズ（\R{全}{ぜん}5\R{巻}{かん}）と、よく\R{使}{つか}う\R{数学}{すう|がく}の\R{巻}{かん}（よく\R{使}{つか}う\R{順}{じゅん}。\R{第}{だい}9\R{巻}{かん}の\R{巻末}{かん|まつ}のカードから）}
\end{bkfigure}


\begin{bkcenter}
\begin{bktable}{colspec={Q[l,m,wd=26mm] Q[c,m,wd=6mm] Q[c,m,wd=6mm] Q[c,m,wd=6mm] Q[c,m,wd=6mm] Q[c,m,wd=6mm] Q[c,m,wd=6mm] Q[c,m,wd=6mm] Q[c,m,wd=6mm] Q[c,m,wd=6mm]}, colsep=1.4mm, rowsep=0pt, cells={font=\small}, row{1}={font=\sffamily\bfseries\small, abovesep=2pt}}
\R{物理}{ぶつ|り}の\R{巻}{かん} & \kgavol{1} & \kgavol{2} & \kgavol{3} & \kgavol{4} & \kgavol{5} & \kgavol{6} & \kgavol{7} & \kgavol{8} & \kgavol{9} \\
1\enspace\R{力学}{りき|がく} & & ● & ● & ● & ● & ● & ● & ● & \\
2\enspace\R{熱学}{ねつ|がく} & & ● & ● & & ● & ● & ● & ● & ● \\
3\enspace\R{波動}{は|どう} & & & ● & ● & ● & ● & ● & ● & ● \\
4\enspace\R{電磁気学}{でん|じ|き|がく} & & & ● & ● & ● & ● & ● & ● & ● \\
5\enspace\R{原子}{げん|し} & & ● & ● & ● & ● & ● & ● & ● & ● \\
\end{bktable}
\bktabcaption{\R{物理}{ぶつ|り}の\R{各巻}{かく|かん}の\R{前}{まえ}に\R{学}{まな}ぶ\R{数学}{すう|がく}の\R{巻}{かん}（●。\R{章}{しょう}ごとの\R{道}{みち}しるべから）}
\end{bkcenter}

\bkfukidashi{sora}{「リンゴは\R{何個}{なん|こ}？」から\R{始}{はじ}めて、\R{工学}{こう|がく}の\R{数学}{すう|がく}の\R{地図}{ち|ず}まで。さあ、\R{自分}{じ|ぶん}の\R{巻}{かん}を\R{開}{ひら}いて、\R{歩}{ある}き\R{始}{はじ}めよう！}

\par\vspace{0.5mm}
\noindent{\sffamily\small\color{bkGray}おうちの\R{人}{ひと}・\R{先生}{せん|せい}には、\R{大人}{おと|な}\R{用}{よう}の『\R{保護者}{ほ|ご|しゃ}・\R{先生}{せん|せい}のための\R{本}{ほん}』があります。}\par
}{\clearpage\typeout{BKPAGE missing guide-p1}}

%% ===================== 奥付・表4 =====================
\makeatletter
\kgaColophon{\bkSeries　ガイドブック}{はじめに\R{読}{よ}む\R{本}{ほん}（\R{子}{こ}ども\R{用}{よう}）}{\R{第}{だい}1\R{版}{はん}（\R{第}{だい}16\R{回}{かい}\R{編集委員会}{へん|しゅう|い|いん|かい}）　2026\R{年}{ねん}10\R{月}{がつ}}%
  {このシリーズは\R{全}{ぜん}9\R{巻}{かん}です。\R{各巻}{かく|かん}の\R{問題}{もん|だい}のくわしい\R{解}{と}き\R{方}{かた}は、これから\R{作}{つく}る\R{別}{べっ}\R{冊}{さつ}の\R{問題集}{もん|だい|しゅう}と\R{指導書}{し|どう|しょ}にのせる\R{予定}{よ|てい}です。}
\bkartpage{表4}{\kgaBackArt{この\R{本}{ほん}でわかること}{%
  \kgaBackItem{197}{1}{このシリーズは、なんの\R{本}{ほん}？}
  \kgaBackItem{188}{2}{9\R{冊}{さつ}の\R{本}{ほん}と、7\R{本}{ほん}の\R{路線}{ろ|せん}}
  \kgaBackItem{179}{3}{1つの\R{章}{しょう}の\R{歩}{ある}き\R{方}{かた}と、★の\R{読}{よ}み\R{方}{かた}}
  \kgaBackItem{170}{4}{どの\R{巻}{かん}から\R{読}{よ}む？}
  \kgaBackItem{161}{5}{\R{勉強}{べん|きょう}のしかたと、\R{巻末}{かん|まつ}の\R{使}{つか}い\R{方}{かた}}}%
  {どの\R{巻}{かん}から\R{読}{よ}む？\\この\R{本}{ほん}で\R{決}{き}めよう}}
\makeatother
\end{document}
